% =====================================================================
% Chapter 7: Eigenvectors, resonator participation and the MAC
% Theory session, dispersion branch (Nov 2026)
% Requires: amsmath, amssymb
% Include in the main file with  % =====================================================================
% Chapter 7: Eigenvectors, resonator participation and the MAC
% Theory session, dispersion branch (Nov 2026)
% Requires: amsmath, amssymb
% Include in the main file with  % =====================================================================
% Chapter 7: Eigenvectors, resonator participation and the MAC
% Theory session, dispersion branch (Nov 2026)
% Requires: amsmath, amssymb
% Include in the main file with  % =====================================================================
% Chapter 7: Eigenvectors, resonator participation and the MAC
% Theory session, dispersion branch (Nov 2026)
% Requires: amsmath, amssymb
% Include in the main file with  \include{chapter_07_eigenvectors_mac}
% =====================================================================

\chapter{Eigenvectors, resonator participation and the MAC}
\label{ch:eigenvectors}

So far only the eigenvalues of the Bloch problem were used. This chapter looks at
the eigenvectors: what they contain, how much of a mode sits in the resonator, how
two modes at neighbouring wavevectors are compared with the Modal Assurance
Criterion (MAC), and which situations a mode tracker has to handle. All FE numbers
are measured on the production cell: $a = 0.1\,\mathrm{m}$, $h = 5\,\mathrm{mm}$,
steel, $10\times 10$ Lagrange degree 2 elements, one resonator at the cell centre
with $f_r = 1200\,\mathrm{Hz}$ and $\mu = 0.25$.

\begin{center}
\begin{tabular}{|l|p{0.68\textwidth}|}
\hline
Symbol & Meaning \\
\hline
$K_r$, $M_r$ & reduced stiffness and mass matrices at one $\mathbf{k}$, $K_r = T^H K T$, $M_r = T^H M T$ \\
$T(\mathbf{k})$ & maps independent DOFs to all DOFs, with phase factors \\
$\lambda = \omega^2$ & eigenvalue of one dot in the diagram \\
$\psi$ & reduced eigenvector of one dot, the column returned by the solver \\
$u$ & last entry of $\psi$: the resonator displacement \\
$m$, $k_s$ & resonator mass and spring stiffness \\
$\lambda_r = k_s/m$ & the resonator's own eigenvalue \\
$\lambda_p$ & bare plate eigenvalue of one plate wave \\
$\mu = m/(\rho h A)$ & resonator mass over plate mass per cell \\
$\langle \psi_1, \psi_2\rangle$ & mass-weighted inner product $\psi_1^H M_r \psi_2$ \\
$H$ & $2\times 2$ standard-form matrix of the two-level model \\
$q$, $v$ & entries of the eigenvector of $H$: plate wave amplitude and $v = \sqrt{m}\,u$ \\
$\bar{H}$, $\Delta$ & mean and half difference of the diagonal entries of $H$ \\
$c$ & off-diagonal entry of $H$ \\
$R$ & $\sqrt{\Delta^2 + c^2}$ \\
$\chi$ & mixing angle of the $2\times 2$ eigenvector \\
$k$ (in tables) & position on the diagonal, $k_x = k_y = k\,\pi/a$ \\
\hline
\end{tabular}
\end{center}

% ---------------------------------------------------------------------
\section{What an eigenvector is at one wavevector}

At each $\mathbf{k}$ the solver returns pairs $(\lambda, \psi)$ with
\begin{equation}
  K_r \psi = \lambda M_r \psi .
\end{equation}
Each dot in the diagram owns one $\psi$. Its entries are the plate unknowns
(deflection and two rotations) at the independent nodes, and its last entry is the
resonator displacement $u$. The resonator has no periodic partner, so $T$ acts on it
as the identity, and the last row of $M_r$ contains only $m$ on the diagonal.

$\psi$ is defined only up to a factor: if $\psi$ solves the equation, so does
$\gamma\psi$ for any complex number $\gamma$. Its scale and its overall complex
phase therefore carry no physical meaning, and everything computed from $\psi$ has
to be independent of both.

% ---------------------------------------------------------------------
\section{Orthogonality and what a degeneracy does}
\label{sec:orthogonality}

Take two solutions at the same $\mathbf{k}$:
\begin{equation}
  K_r \psi_1 = \lambda_1 M_r \psi_1, \qquad K_r \psi_2 = \lambda_2 M_r \psi_2 .
\end{equation}
Multiply the first from the left by $\psi_2^H$:
\begin{equation}
  \psi_2^H K_r \psi_1 = \lambda_1\, \psi_2^H M_r \psi_1 .
\end{equation}
Take the conjugate transpose of the second. $K_r$ and $M_r$ are Hermitian and
$\lambda_2$ is real, so $\psi_2^H K_r = \lambda_2\, \psi_2^H M_r$. Multiply from the
right by $\psi_1$:
\begin{equation}
  \psi_2^H K_r \psi_1 = \lambda_2\, \psi_2^H M_r \psi_1 .
\end{equation}
The left sides are equal. Subtracting gives
\begin{equation}
  (\lambda_1 - \lambda_2)\, \psi_2^H M_r \psi_1 = 0 .
\end{equation}
Different eigenvalues therefore mean $\psi_2^H M_r \psi_1 = 0$: the vectors are
orthogonal when weighted with $M_r$. That is why the mass matrix appears in
everything below.

If $\lambda_1 = \lambda_2 = \lambda$, the equation says nothing. Then every
combination is also a solution,
\begin{equation}
  K_r(\alpha_1\psi_1 + \alpha_2\psi_2) = \lambda M_r(\alpha_1\psi_1 + \alpha_2\psi_2),
\end{equation}
so at a degenerate point the solver returns \emph{some arbitrary basis} of the
group. The individual vectors of the group have no meaning of their own; only the
space they span does.

In this cell, degenerate groups appear where several bare plate waves share one
frequency. One combination couples to the resonator and leaves the group; the rest
do not see it. For example, $M$ has four plate waves at $2378\,\mathrm{Hz}$: one
coupled combination moves away and three modes remain at $2379\,\mathrm{Hz}$.

% ---------------------------------------------------------------------
\section{Plate share in the FE model}

The kinetic energy of a mode is $\tfrac12\omega^2\,\psi^H M_r\psi$, and the part in
the resonator is $\tfrac12\omega^2\, m|u|^2$. Hence
\begin{equation}
  \text{plate share} = 1 - \frac{m|u|^2}{\psi^H M_r \psi} .
\end{equation}
Scale and phase of $\psi$ cancel, as required.

\begin{center}
\begin{tabular}{|l|p{0.72\textwidth}|}
\hline
Point & Frequency in Hz: plate share (FE) \\
\hline
$\Gamma$ & $0$: $80\,\%$;\quad $1274$: $20\,\%$;\quad $4730$: $100\,\%$ (three modes);\quad $4883$: $100\,\%$ \\
$X$ & $836$: $31\,\%$;\quad $1195$: $100\,\%$;\quad $1653$: $69\,\%$;\quad $5881$: $100\,\%$ (three modes) \\
$M$ & $1031$: $5\,\%$;\quad $2379$: $100\,\%$ (three modes);\quad $2712$: $95\,\%$;\quad $11567$: $100\,\%$ \\
\hline
\end{tabular}
\end{center}

The simple model of Chapter~\ref{ch:local_resonance} gave $34\,\%$ and $66\,\%$ at
$X$; the shares at $\Gamma$ and $M$ agree with it.

\paragraph{The sum rule for the resonator share.}
Collect the eigenvectors of all bands at one $\mathbf{k}$ as columns of a matrix
$\Psi$, normalized so that $\Psi^H M_r \Psi = I$. Then $\Psi^{-1} = \Psi^H M_r$, so
\begin{equation}
  \Psi\,\Psi^H M_r = I \quad\Rightarrow\quad \Psi\,\Psi^H = M_r^{-1} .
\end{equation}
The resonator has no mass coupling to the plate, so the last diagonal entry of
$M_r^{-1}$ is $1/m$. The last diagonal entry of $\Psi\Psi^H$ is $\sum_n |u_n|^2$,
with $n$ running over all bands. Therefore
\begin{equation}
  \sum_n m|u_n|^2 = 1 .
\end{equation}
With this normalization $m|u_n|^2$ is the resonator share of band $n$, so
\emph{the resonator shares of all bands at one $\mathbf{k}$ add up to exactly
$100\,\%$}. The FE values at the three symmetry points show it: $20 + 80$ at
$\Gamma$, $69 + 31$ at $X$, $95 + 5$ at $M$. The plate shares of bands 1 and 2 alone
need not add up to $100\,\%$: along $M$--$\Gamma$ their sum drifts from $100\,\%$ to
about $104\,\%$ close to $M$, because other bands take part there.

% ---------------------------------------------------------------------
\section{The Modal Assurance Criterion}

\subsection{Definition and properties}
We need a number that says how alike two eigenvectors are and that ignores scale
and phase. With $\langle \psi_1, \psi_2\rangle = \psi_1^H M_r \psi_2$,
\begin{equation}
  \mathrm{MAC}(\psi_1, \psi_2) = \frac{|\langle\psi_1, \psi_2\rangle|^2}{\langle\psi_1, \psi_1\rangle\,\langle\psi_2, \psi_2\rangle} .
\end{equation}
\begin{itemize}
  \item Scale and phase drop out: replacing $\psi_1$ by $\gamma\psi_1$ multiplies
        numerator and denominator both by $|\gamma|^2$.
  \item By the Cauchy--Schwarz inequality it lies between 0 and 1. It is the squared
        cosine of the angle between the two vectors.
  \item $\mathrm{MAC} = 1$ for the same shape.
  \item $\mathrm{MAC} = 0$ for $M_r$-orthogonal vectors, which by
        Section~\ref{sec:orthogonality} holds for two different bands at the same
        $\mathbf{k}$.
\end{itemize}
Two vectors at neighbouring $\mathbf{k}$ come with two slightly different $M_r$. The
measurements below use $M_r$ of the first of the two points; the difference
vanishes as the step shrinks.

\subsection{How the step size enters}
Normalize $\langle\psi,\psi\rangle = 1$ and write the vector at the next
wavevector as
\begin{equation}
  \psi(\mathbf{k} + \delta\mathbf{k}) \approx \psi + \delta\,(\alpha\,\psi + p),
  \qquad \langle\psi, p\rangle = 0,
\end{equation}
with $\delta$ the step length, $\alpha$ a number and $p$ the part of the change
perpendicular to $\psi$. Then
\begin{align}
  \langle\psi, \psi(\mathbf{k} + \delta\mathbf{k})\rangle &= 1 + \delta\alpha, \\
  \langle\psi(\mathbf{k} + \delta\mathbf{k}), \psi(\mathbf{k} + \delta\mathbf{k})\rangle &= |1 + \delta\alpha|^2 + \delta^2\langle p, p\rangle ,
\end{align}
and so
\begin{equation}
  \mathrm{MAC} = \frac{|1 + \delta\alpha|^2}{|1 + \delta\alpha|^2 + \delta^2\langle p,p\rangle}
  \approx 1 - \delta^2\langle p, p\rangle .
\end{equation}
The part along $\psi$ only changes scale and phase and drops out. So
$1 - \mathrm{MAC} \propto \delta^2$, and halving the step should reduce it by a
factor of 4. Measured for band 1 at its worst step over the whole path (the point
number refers to the segments $\Gamma$--$X$ and $X$--$M$; $M$--$\Gamma$ gets
proportionally more):

\begin{center}
\begin{tabular}{|l|l|l|l|}
\hline
Points/segment & Wavevectors & Lowest MAC, band 1 & $1 - \mathrm{MAC}$ \\
\hline
20 & 69 & 0.9622 & 0.0378 \\
40 & 138 & 0.9898 & 0.0102 \\
80 & 274 & 0.9974 & 0.0026 \\
\hline
\end{tabular}
\end{center}

The ratios are $3.7$ and $3.9$, as expected. At 20 points per segment this worst
step lies just before $X$ on $\Gamma$--$X$, not at the resonator crossing on
$M$--$\Gamma$.

\subsection{Why only neighbours are compared}
Band 1 at $\Gamma$ against band 1 at $M$ has $\mathrm{MAC} = 0.35$, because the mode
legitimately changes from $80\,\%$ plate to $5\,\%$ plate. A low value over a long
distance means nothing. Between neighbouring points on $M$--$\Gamma$ the same band
never drops below $0.9974$ at 40 points per segment, even though it changes its
character completely along that path.

\subsection{Which vector and which weighting}
Lowest neighbouring MAC of band 1 at 40 points per segment:

\begin{center}
\begin{tabular}{|l|l|}
\hline
Variant & Lowest value \\
\hline
Full vector, $M_r$-weighted & 0.990 \\
Plate entries only, $M_r$-weighted & 0.969 \\
Full vector, unweighted & 0.619 \\
\hline
\end{tabular}
\end{center}

The MAC normalizes, so what is left of the plate motion in a resonator-dominated mode
still gives a usable comparison, only slightly worse. The clearly bad variant is
leaving out $M_r$: it counts every entry equally, whether deflection, rotation or
resonator, and falls to $0.62$ on a band that is perfectly smooth. Stage 5 therefore
uses the full vector including $u$, weighted with $M_r$.

% ---------------------------------------------------------------------
\section{How fast a vector turns at an avoided crossing}
\label{sec:mixing_angle}

This section takes the $2\times 2$ matrix of the two-level model apart, to see how
the eigenvector turns through an avoided crossing.

\subsection{The matrix we start from}
For one plate wave and one resonator, in mass-normalized coordinates, the eigenvector
is $(q, v)$ with $q$ the plate wave amplitude and $v = \sqrt{m}\,u$, and
Chapter~\ref{ch:local_resonance} gave
\begin{equation}
  H = \begin{pmatrix} H_{11} & c \\ c & H_{22} \end{pmatrix}
    = \begin{pmatrix} \lambda_p + \mu\lambda_r & -\sqrt{\mu}\,\lambda_r \\ -\sqrt{\mu}\,\lambda_r & \lambda_r \end{pmatrix} .
\end{equation}

\subsection{Mean and half difference}
Define the mean and the half difference of the diagonal,
\begin{equation}
  \bar{H} = \frac{H_{11} + H_{22}}{2}, \qquad \Delta = \frac{H_{11} - H_{22}}{2} .
\end{equation}
Adding and subtracting these gives $H_{11} = \bar{H} + \Delta$ and
$H_{22} = \bar{H} - \Delta$, so
\begin{equation}
  H = \bar{H}\, I + B, \qquad B = \begin{pmatrix} \Delta & c \\ c & -\Delta \end{pmatrix} .
\end{equation}
Inserting the entries:

\begin{center}
\begin{tabular}{|l|l|l|}
\hline
Symbol & In terms of $H$ & For plate wave and resonator \\
\hline
$\bar{H}$ & $(H_{11} + H_{22})/2$ & $\tfrac12\left[\lambda_p + \lambda_r(1+\mu)\right]$ \\
$\Delta$ & $(H_{11} - H_{22})/2$ & $\tfrac12\left[\lambda_p - \lambda_r(1-\mu)\right]$ \\
$c$ & off-diagonal entry & $-\sqrt{\mu}\,\lambda_r$ \\
\hline
\end{tabular}
\end{center}

\paragraph{Why split it.} If $Bz = rz$ for a vector $z$, then
$Hz = \bar{H}z + Bz = (\bar{H} + r)z$. The part $\bar{H}I$ only shifts the
eigenvalues. \emph{The eigenvectors depend on $B$ alone}, so only $\Delta$ and $c$
matter. $\Delta = 0$ at $\lambda_p = \lambda_r(1-\mu)$, where the diagonal entries
are equal: the centre of the avoided crossing.

\subsection{Eigenvalues}
\begin{align}
  \det(B - rI) &= (\Delta - r)(-\Delta - r) - c^2 = r^2 - \Delta^2 - c^2 = 0, \\
  \Rightarrow\quad r &= \pm R, \qquad R = \sqrt{\Delta^2 + c^2} .
\end{align}
The eigenvalues of $H$ are $\bar{H} \pm R$. This is the familiar two-level formula
\begin{equation}
  \lambda = \frac{H_{11} + H_{22}}{2} \pm \sqrt{\left(\frac{H_{11} - H_{22}}{2}\right)^2 + |c|^2} .
\end{equation}
Their sum is $2\bar{H} = \lambda_p + \lambda_r(1+\mu)$ and their product is
\begin{equation}
  \bar{H}^2 - R^2 = \bar{H}^2 - \Delta^2 - c^2 = H_{11}H_{22} - c^2
  = (\lambda_p + \mu\lambda_r)\lambda_r - \mu\lambda_r^2 = \lambda_p\lambda_r .
\end{equation}

\subsection{Eigenvectors and the mixing angle}
$(\Delta, c)$ is a point in a plane at distance $R$ from the origin. Call its polar
angle $2\chi$:
\begin{equation}
  \Delta = R\cos 2\chi, \qquad c = R\sin 2\chi, \qquad \tan 2\chi = \frac{c}{\Delta} .
\end{equation}
This is always possible because $\Delta^2 + c^2 = R^2$. Then
\begin{equation}
  B = R \begin{pmatrix} \cos 2\chi & \sin 2\chi \\ \sin 2\chi & -\cos 2\chi \end{pmatrix} .
\end{equation}
Apply $B$ to $(\cos\chi, \sin\chi)$ and use the addition theorems for
$\cos(2\chi - \chi)$ and $\sin(2\chi - \chi)$:
\begin{align}
  \text{row 1:} \quad & R(\cos 2\chi\cos\chi + \sin 2\chi\sin\chi) = R\cos(2\chi - \chi) = R\cos\chi, \\
  \text{row 2:} \quad & R(\sin 2\chi\cos\chi - \cos 2\chi\sin\chi) = R\sin(2\chi - \chi) = R\sin\chi .
\end{align}
So $(\cos\chi, \sin\chi)$ belongs to $+R$. Applying $B$ to $(-\sin\chi, \cos\chi)$
in the same way:
\begin{align}
  \text{row 1:} \quad & R(-\cos 2\chi\sin\chi + \sin 2\chi\cos\chi) = R\sin\chi = -R\,(-\sin\chi), \\
  \text{row 2:} \quad & R(-\sin 2\chi\sin\chi - \cos 2\chi\cos\chi) = -R\cos\chi ,
\end{align}
so it belongs to $-R$. In summary,
\begin{equation}
  \text{upper band: } (q, v) = (\cos\chi, \sin\chi), \qquad
  \text{lower band: } (q, v) = (-\sin\chi, \cos\chi) .
\end{equation}
The plate share is $|q|^2$. With $\cos^2\chi = (1 + \cos 2\chi)/2$ and
$\cos 2\chi = \Delta/R$,
\begin{equation}
  \boxed{\;\begin{aligned}
  \text{upper band: plate share} &= \tfrac12\left(1 + \frac{\Delta}{R}\right), \\
  \text{lower band: plate share} &= \tfrac12\left(1 - \frac{\Delta}{R}\right)
  \end{aligned}\;}
\end{equation}
Because $c < 0$ here, $\chi$ is negative. That only sets the relative sign of $q$ and
$v$; the shares do not depend on it.

\subsection{One worked point: $\Gamma$}
At $\Gamma$, $\lambda_p = 0$, and with $\mu = 0.25$:
\begin{align}
  H_{11} &= 0.25\,\lambda_r, \quad H_{22} = \lambda_r, \quad c = -0.5\,\lambda_r, \\
  \bar{H} &= 0.625\,\lambda_r, \quad \Delta = -0.375\,\lambda_r, \\
  R &= \sqrt{0.1406 + 0.25}\,\lambda_r = 0.625\,\lambda_r .
\end{align}
The eigenvalues are $\bar{H} \pm R = 0$ and $1.25\,\lambda_r$, the known $0$ and
$\lambda_r(1+\mu)$. With $\Delta/R = -0.6$, the lower band has plate share
$\tfrac12(1 + 0.6) = 80\,\%$ and the upper band $\tfrac12(1 - 0.6) = 20\,\%$.
Further, $2\chi = -126.9^\circ$, so $\chi = -63.4^\circ$.

This agrees with the ratio formula of Chapter~\ref{ch:local_resonance},
$q/v = (\lambda_r - \lambda)/(\sqrt{\mu}\,\lambda_r)$: for the upper band it gives
$-0.25/0.5 = -0.5$, and $\cot(-63.4^\circ) = -0.5$.

It also explains why band 1 starts at $80\,\%$ and not at $100\,\%$. Since $\lambda_p$
cannot go below 0, $\Delta$ cannot go below $-0.375\,\lambda_r$, and the far end of
the turn is never reached.

\subsection{Along $M$--$\Gamma$: model against FE}
Plate shares in percent, for bands 1 and 2 (B1, B2), one-wave model (mod.)\ against
the FE cell:

\begin{center}
\begin{tabular}{|l|l|l|l|l|l|l|l|}
\hline
$k$ & $\lambda_p/\lambda_r$ & $\Delta/\lambda_r$ & $\chi$ & B1 mod. & B1 FE & B2 mod. & B2 FE \\
\hline
0.018 & 0.000 & $-0.375$ & $-63.4^\circ$ & 80.0 & 80.0 & 20.0 & 20.4 \\
0.509 & 0.267 & $-0.242$ & $-57.9^\circ$ & 71.8 & 70.3 & 28.2 & 31.2 \\
0.649 & 0.705 & $-0.023$ & $-46.3^\circ$ & 52.3 & 46.3 & 47.7 & 56.1 \\
0.719 & 1.060 & $0.155$ & $-36.4^\circ$ & 35.2 & 27.7 & 64.8 & 75.2 \\
0.789 & 1.536 & $0.393$ & $-25.9^\circ$ & 19.1 & 14.2 & 80.9 & 89.2 \\
0.930 & 2.943 & $1.096$ & $-12.3^\circ$ & 4.5 & 5.6 & 95.5 & 98.9 \\
1.000 & 3.928 & $1.589$ & $-8.7^\circ$ & 2.3 & 5.1 & 97.7 & 100.0 \\
\hline
\end{tabular}
\end{center}

The model has $\Delta = 0$ at $k \approx 0.667$, and the FE band 1 passes $50\,\%$
shortly before $k = 0.649$. The agreement fades towards $M$, as it should: at $M$
four plate waves take part, and this matrix holds only one. The band frequencies show
the same (band 1 at $M$: $1153\,\mathrm{Hz}$ in the model, $1031\,\mathrm{Hz}$ in
FE). With these values $\chi$ runs from $-63.4^\circ$ at $\Gamma$ to $-8.7^\circ$ at
$M$, about $55^\circ$ in total, not the idealized full turn of $90^\circ$.

\subsection{From the angle to the MAC}
Two normalized vectors at angles $\chi_1$ and $\chi_2$ have the inner product
$\cos\chi_1\cos\chi_2 + \sin\chi_1\sin\chi_2 = \cos(\chi_1 - \chi_2)$, so
\begin{equation}
  \mathrm{MAC} = \cos^2(\chi_1 - \chi_2) .
\end{equation}
To see how fast $\chi$ turns, differentiate $2\chi = \arctan(c/\Delta)$ with respect
to $\Delta$:
\begin{equation}
  2\,\frac{\mathrm{d}\chi}{\mathrm{d}\Delta}
  = \frac{1}{1 + c^2/\Delta^2}\cdot\left(-\frac{c}{\Delta^2}\right)
  = -\frac{c}{\Delta^2 + c^2}
  = -\frac{c}{R^2} .
\end{equation}
The turn is fastest at $\Delta = 0$, where $R = |c|$ and
$|\mathrm{d}\chi/\mathrm{d}\Delta| = 1/(2|c|)$. Since $\Delta$ contains $\lambda_p/2$,
\begin{equation}
  \left|\frac{\mathrm{d}\chi}{\mathrm{d}\lambda_p}\right|_{\Delta = 0}
  = \frac{1}{4|c|} = \frac{1}{4\sqrt{\mu}\,\lambda_r} .
\end{equation}
Half of the full $90^\circ$ turn happens while $|\Delta| < |c|$: at $\Delta = |c|$ one
has $2\chi = -45^\circ$, at $\Delta = -|c|$ one has $2\chi = -135^\circ$, so $\chi$
changes by $45^\circ$. In terms of $\lambda_p$ this window has the width
$4\sqrt{\mu}\,\lambda_r$: $2.0\,\lambda_r$ for $\mu = 0.25$, $0.4\,\lambda_r$ for
$\mu = 0.01$. A lighter resonator has a narrower crossing and needs smaller steps.

Measured for band 1 on $M$--$\Gamma$, with the spacing $0.0175\,\pi/a$ of the 40-point
path (57 steps on this segment):

\begin{center}
\begin{tabular}{|l|l|l|}
\hline
 & Lowest MAC on $M$--$\Gamma$ & at $k$ \\
\hline
Model, $\cos^2(\chi_1 - \chi_2)$ & 0.99778 & 0.737 \\
FE & 0.99736 & 0.702 \\
\hline
\end{tabular}
\end{center}

The minimum lies a little past $\Delta = 0$ because $\lambda_p$ also changes faster
with $k$ as $k$ grows. The change from $69\,\%$ to $24\,\%$ plate takes 12 steps, and
the change from $80\,\%$ to $5\,\%$ takes practically all 57.

\paragraph{The baseline of the FE value.}
Close to $\Gamma$ the FE MAC stays at $0.99961$ while the model gives $1.00000$. The
model uses ``the plate wave at this $\mathbf{k}$'' as a basis vector, so it cannot see
that the pattern $e^{i\mathbf{k}\cdot\mathbf{x}}$ itself changes between neighbouring
points. One step changes the phase at $\mathbf{x}$ by $\theta(\mathbf{x}) =
\delta\mathbf{k}\cdot\mathbf{x}$, and $1 - \mathrm{MAC}$ is approximately the variance
of $\theta$ over the mass of the mode. With $\delta k_x = \delta k_y = 0.0175\,\pi/a$,
a variance of $a^2/12$ of each coordinate over the cell, and $80\,\%$ of the mass
motion in the plate (the resonator sits at the mean position):
\begin{equation}
  1 - \mathrm{MAC} \approx 0.8\cdot (0.0175\,\pi)^2\cdot\frac{2}{12} \approx 0.0004 ,
\end{equation}
which matches the measured $1 - 0.99961 = 0.00039$.

\subsection{The other two-level matrix: two plate waves and a rigid mass}
Chapter~\ref{ch:two_level} used the same structure with
$H_{11} = \lambda_1^0 - \mu\lambda$, $H_{22} = \lambda_2^0 - \mu\lambda$ and
$c = -\mu\lambda$. Then
\begin{equation}
  \bar{H} = \frac{\lambda_1^0 + \lambda_2^0}{2} - \mu\lambda, \qquad
  \Delta = \frac{\lambda_1^0 - \lambda_2^0}{2},
\end{equation}
because the mass loading cancels in the difference. At $X$, $\Delta = 0$, so
$2\chi = -90^\circ$ and $\chi = -45^\circ$. That gives the eigenvectors
$(1, -1)/\sqrt{2}$ and $(1, 1)/\sqrt{2}$ of that chapter, with $a_H = \bar{H}$ and
$b_H = c$.

% ---------------------------------------------------------------------
\section{The three situations a tracker meets}

Lowest neighbouring MAC between \emph{sorted} bands along the whole path:

\begin{center}
\begin{tabular}{|l|l|l|l|p{0.32\textwidth}|}
\hline
Band & 20 points & 40 points & 80 points & Where \\
\hline
1 & 0.9622 & 0.9898 & 0.9974 & smooth band \\
2 & 0.2301 & 0.2325 & 0.2322 & next to $M$, three modes at $2379\,\mathrm{Hz}$ \\
3 & 0.1755 & 0.1954 & 0.2117 & next to $\Gamma$, three modes at $4730\,\mathrm{Hz}$ \\
4 & 0.0601 & 0.0610 & 0.0615 & next to $X$, three modes at $5881\,\mathrm{Hz}$ \\
5 & 0.0000 & 0.0000 & 0.0000 & near $k \approx 0.33$ on $M$--$\Gamma$, about $6.5\,\mathrm{kHz}$ \\
\hline
\end{tabular}
\end{center}

\paragraph{1. Smooth continuation.} Band 1: the MAC is close to 1 and improves with a
finer step.

\paragraph{2. A degenerate point.} Bands 2 to 4: the value barely changes with a finer
step. That is the signature of a degeneracy and not of a coarse step. At the
degenerate point the solver picks an arbitrary basis of the group
(Section~\ref{sec:orthogonality}), while just beside it the modes are distinct. Near
$M$, MAC of the three neighbouring modes against the three modes at $M$ (40 points per
segment):

\begin{center}
\begin{tabular}{|l|l|l|l|l|}
\hline
Before $M$ & $M$, mode 1 & $M$, mode 2 & $M$, mode 3 & Row sum \\
\hline
$2321\,\mathrm{Hz}$ & 0.497 & 0.045 & 0.457 & 1.000 \\
$2372\,\mathrm{Hz}$ & 0.017 & 0.882 & 0.072 & 0.971 \\
$2439\,\mathrm{Hz}$ & 0.485 & 0.047 & 0.468 & 1.000 \\
\hline
\end{tabular}
\end{center}

No single entry tells the continuation, but the row sums are close to 1. The reason:
let $e_1, \dots, e_g$ be an $M_r$-orthonormal basis of the group and $\psi$ a vector
with $\langle\psi,\psi\rangle = 1$. Then
\begin{equation}
  \sum_{j=1}^{g} \mathrm{MAC}(\psi, e_j) = \sum_{j=1}^{g} |\langle e_j, \psi\rangle|^2,
\end{equation}
which by Pythagoras is the squared length of the part of $\psi$ that lies inside the
group. It equals 1 exactly if $\psi$ lies entirely in the group, and it does not depend
on which basis the solver chose. The neighbouring mode sits at a slightly different
$\mathbf{k}$, so the sum is only approximately 1. Degenerate groups therefore have to
be compared as a whole, and all members of the group must have been computed.

\paragraph{3. A true crossing.} Two sorted bands swap identities: the sorted MAC is
about 0 and the MAC with the other band is about 1. In the sweep, bands 5 and 6 near
$k \approx 0.33$ on $M$--$\Gamma$ (about $6.5\,\mathrm{kHz}$, above the $5\,\mathrm{kHz}$
range) show a MAC of $0.0000$. That is consistent with two modes of opposite symmetry
about the diagonal, which have $c = 0$; the symmetry itself was not checked.

% ---------------------------------------------------------------------
\section{Tracking and gap finding are different tasks}

\emph{Tracking} follows a mode's identity from one $\mathbf{k}$ to the next, using the
MAC between neighbours. It answers ``which curve is this dot on?''

\emph{Gap finding} asks whether any dot lies in a frequency interval. It needs only the
frequencies sorted at each $\mathbf{k}$: the gap between sorted bands $n$ and $n+1$
runs from the maximum of band $n$ to the minimum of band $n+1$.

A tracked band may cross another one, so the $n$-th tracked band is not the $n$-th
sorted band. Suppose tracked curves $A$ and $B$ cross and a third band $C$ lies above
both. The sorted band below $C$ is the upper envelope of $A$ and $B$ at each
$\mathbf{k}$, so the correct lower gap edge is the maximum of that envelope. Taking the
curve $B$ alone when $A$ rises higher later on the path gives a value that is too low:
the reported gap starts too low and contains dots of $A$. Stage 5 therefore keeps plain
sorting available.

% ---------------------------------------------------------------------
\section{Check questions with answers}

\paragraph{Question 1.} \emph{Band 3 near $\Gamma$ has a neighbouring MAC of about 0.2,
and it stays there when the number of $\mathbf{k}$-points is doubled. What do you
conclude, and what would you do differently for that step?}

\emph{Answer.} Band 3 meets two other bands at $\Gamma$, and the solver returns an
arbitrary basis of that group. The MAC has to be summed over the whole group, which
then gives approximately 1. A degeneracy has to be told apart from a true crossing
between two sampled points: there the continuation is well defined, the MAC with the
partner band is about 1 and the sorted labels simply swap. Here the MAC is about 0.2,
neither 0 nor 1, and refining the step does not change it, which points to a degenerate
point where individual vectors mean nothing. The sum is only approximately 1 because the
neighbouring vector sits at a slightly different $\mathbf{k}$ (measured: $0.97$ to
$1.00$), and it needs every member of the group, which is why the solver uses guard
eigenpairs. For that step: match the group as a whole instead of one vector to one
vector. At the degenerate dot all members have the same frequency, so the plotted
curve does not depend on which label goes where.

\paragraph{Question 2.} \emph{Band 1 changes from $80\,\%$ plate to $5\,\%$ plate along
$M$--$\Gamma$, yet the neighbouring MAC stays close to 1 there. Why is this not a
contradiction, and which change in the unit cell would make this crossing harder to
track?}

\emph{Answer.} The change is gradual, spread over essentially all 57 steps of the
segment ($69\,\%$ to $24\,\%$ alone takes 12 steps), and each step only compares two
nearly identical shapes. A lighter resonator makes the crossing narrower, because the
coupling is $c = -\sqrt{\mu}\,\lambda_r$; the turn then happens within fewer steps, and
at the same step size the MAC dips further from 1.

\paragraph{Question 3.} \emph{Two tracked bands overlap in frequency, with a true crossing
between them. Could there still be a complete gap between sorted bands $n$ and $n+1$ at
that crossing frequency, and which data decides it?}

\emph{Answer.} No. At a true crossing the sorted bands $n$ and $n+1$ share a dot at the
crossing frequency $f_c$, so $\max(\text{band } n) \ge f_c \ge \min(\text{band } n+1)$
and there is no gap between that pair. The data to decide this is the sorted frequencies
at every $\mathbf{k}$; no eigenvectors are needed. The point of the question is the
reverse danger: using tracked bands for gap edges can report a gap that contains dots of
another band, as shown in the previous section.

\paragraph{Follow-up: how does the matrix $H$ of the two-level model translate into the
notation with mean, half difference and mixing angle?}
\emph{Answer.} This is worked out in Section~\ref{sec:mixing_angle}: $H_{11}$ and
$H_{22}$ are replaced by $\bar{H} \pm \Delta$, the off-diagonal $c$ stays the same, and
the eigenvectors depend only on $c/\Delta$.

% =====================================================================

\chapter{Eigenvectors, resonator participation and the MAC}
\label{ch:eigenvectors}

So far only the eigenvalues of the Bloch problem were used. This chapter looks at
the eigenvectors: what they contain, how much of a mode sits in the resonator, how
two modes at neighbouring wavevectors are compared with the Modal Assurance
Criterion (MAC), and which situations a mode tracker has to handle. All FE numbers
are measured on the production cell: $a = 0.1\,\mathrm{m}$, $h = 5\,\mathrm{mm}$,
steel, $10\times 10$ Lagrange degree 2 elements, one resonator at the cell centre
with $f_r = 1200\,\mathrm{Hz}$ and $\mu = 0.25$.

\begin{center}
\begin{tabular}{|l|p{0.68\textwidth}|}
\hline
Symbol & Meaning \\
\hline
$K_r$, $M_r$ & reduced stiffness and mass matrices at one $\mathbf{k}$, $K_r = T^H K T$, $M_r = T^H M T$ \\
$T(\mathbf{k})$ & maps independent DOFs to all DOFs, with phase factors \\
$\lambda = \omega^2$ & eigenvalue of one dot in the diagram \\
$\psi$ & reduced eigenvector of one dot, the column returned by the solver \\
$u$ & last entry of $\psi$: the resonator displacement \\
$m$, $k_s$ & resonator mass and spring stiffness \\
$\lambda_r = k_s/m$ & the resonator's own eigenvalue \\
$\lambda_p$ & bare plate eigenvalue of one plate wave \\
$\mu = m/(\rho h A)$ & resonator mass over plate mass per cell \\
$\langle \psi_1, \psi_2\rangle$ & mass-weighted inner product $\psi_1^H M_r \psi_2$ \\
$H$ & $2\times 2$ standard-form matrix of the two-level model \\
$q$, $v$ & entries of the eigenvector of $H$: plate wave amplitude and $v = \sqrt{m}\,u$ \\
$\bar{H}$, $\Delta$ & mean and half difference of the diagonal entries of $H$ \\
$c$ & off-diagonal entry of $H$ \\
$R$ & $\sqrt{\Delta^2 + c^2}$ \\
$\chi$ & mixing angle of the $2\times 2$ eigenvector \\
$k$ (in tables) & position on the diagonal, $k_x = k_y = k\,\pi/a$ \\
\hline
\end{tabular}
\end{center}

% ---------------------------------------------------------------------
\section{What an eigenvector is at one wavevector}

At each $\mathbf{k}$ the solver returns pairs $(\lambda, \psi)$ with
\begin{equation}
  K_r \psi = \lambda M_r \psi .
\end{equation}
Each dot in the diagram owns one $\psi$. Its entries are the plate unknowns
(deflection and two rotations) at the independent nodes, and its last entry is the
resonator displacement $u$. The resonator has no periodic partner, so $T$ acts on it
as the identity, and the last row of $M_r$ contains only $m$ on the diagonal.

$\psi$ is defined only up to a factor: if $\psi$ solves the equation, so does
$\gamma\psi$ for any complex number $\gamma$. Its scale and its overall complex
phase therefore carry no physical meaning, and everything computed from $\psi$ has
to be independent of both.

% ---------------------------------------------------------------------
\section{Orthogonality and what a degeneracy does}
\label{sec:orthogonality}

Take two solutions at the same $\mathbf{k}$:
\begin{equation}
  K_r \psi_1 = \lambda_1 M_r \psi_1, \qquad K_r \psi_2 = \lambda_2 M_r \psi_2 .
\end{equation}
Multiply the first from the left by $\psi_2^H$:
\begin{equation}
  \psi_2^H K_r \psi_1 = \lambda_1\, \psi_2^H M_r \psi_1 .
\end{equation}
Take the conjugate transpose of the second. $K_r$ and $M_r$ are Hermitian and
$\lambda_2$ is real, so $\psi_2^H K_r = \lambda_2\, \psi_2^H M_r$. Multiply from the
right by $\psi_1$:
\begin{equation}
  \psi_2^H K_r \psi_1 = \lambda_2\, \psi_2^H M_r \psi_1 .
\end{equation}
The left sides are equal. Subtracting gives
\begin{equation}
  (\lambda_1 - \lambda_2)\, \psi_2^H M_r \psi_1 = 0 .
\end{equation}
Different eigenvalues therefore mean $\psi_2^H M_r \psi_1 = 0$: the vectors are
orthogonal when weighted with $M_r$. That is why the mass matrix appears in
everything below.

If $\lambda_1 = \lambda_2 = \lambda$, the equation says nothing. Then every
combination is also a solution,
\begin{equation}
  K_r(\alpha_1\psi_1 + \alpha_2\psi_2) = \lambda M_r(\alpha_1\psi_1 + \alpha_2\psi_2),
\end{equation}
so at a degenerate point the solver returns \emph{some arbitrary basis} of the
group. The individual vectors of the group have no meaning of their own; only the
space they span does.

In this cell, degenerate groups appear where several bare plate waves share one
frequency. One combination couples to the resonator and leaves the group; the rest
do not see it. For example, $M$ has four plate waves at $2378\,\mathrm{Hz}$: one
coupled combination moves away and three modes remain at $2379\,\mathrm{Hz}$.

% ---------------------------------------------------------------------
\section{Plate share in the FE model}

The kinetic energy of a mode is $\tfrac12\omega^2\,\psi^H M_r\psi$, and the part in
the resonator is $\tfrac12\omega^2\, m|u|^2$. Hence
\begin{equation}
  \text{plate share} = 1 - \frac{m|u|^2}{\psi^H M_r \psi} .
\end{equation}
Scale and phase of $\psi$ cancel, as required.

\begin{center}
\begin{tabular}{|l|p{0.72\textwidth}|}
\hline
Point & Frequency in Hz: plate share (FE) \\
\hline
$\Gamma$ & $0$: $80\,\%$;\quad $1274$: $20\,\%$;\quad $4730$: $100\,\%$ (three modes);\quad $4883$: $100\,\%$ \\
$X$ & $836$: $31\,\%$;\quad $1195$: $100\,\%$;\quad $1653$: $69\,\%$;\quad $5881$: $100\,\%$ (three modes) \\
$M$ & $1031$: $5\,\%$;\quad $2379$: $100\,\%$ (three modes);\quad $2712$: $95\,\%$;\quad $11567$: $100\,\%$ \\
\hline
\end{tabular}
\end{center}

The simple model of Chapter~\ref{ch:local_resonance} gave $34\,\%$ and $66\,\%$ at
$X$; the shares at $\Gamma$ and $M$ agree with it.

\paragraph{The sum rule for the resonator share.}
Collect the eigenvectors of all bands at one $\mathbf{k}$ as columns of a matrix
$\Psi$, normalized so that $\Psi^H M_r \Psi = I$. Then $\Psi^{-1} = \Psi^H M_r$, so
\begin{equation}
  \Psi\,\Psi^H M_r = I \quad\Rightarrow\quad \Psi\,\Psi^H = M_r^{-1} .
\end{equation}
The resonator has no mass coupling to the plate, so the last diagonal entry of
$M_r^{-1}$ is $1/m$. The last diagonal entry of $\Psi\Psi^H$ is $\sum_n |u_n|^2$,
with $n$ running over all bands. Therefore
\begin{equation}
  \sum_n m|u_n|^2 = 1 .
\end{equation}
With this normalization $m|u_n|^2$ is the resonator share of band $n$, so
\emph{the resonator shares of all bands at one $\mathbf{k}$ add up to exactly
$100\,\%$}. The FE values at the three symmetry points show it: $20 + 80$ at
$\Gamma$, $69 + 31$ at $X$, $95 + 5$ at $M$. The plate shares of bands 1 and 2 alone
need not add up to $100\,\%$: along $M$--$\Gamma$ their sum drifts from $100\,\%$ to
about $104\,\%$ close to $M$, because other bands take part there.

% ---------------------------------------------------------------------
\section{The Modal Assurance Criterion}

\subsection{Definition and properties}
We need a number that says how alike two eigenvectors are and that ignores scale
and phase. With $\langle \psi_1, \psi_2\rangle = \psi_1^H M_r \psi_2$,
\begin{equation}
  \mathrm{MAC}(\psi_1, \psi_2) = \frac{|\langle\psi_1, \psi_2\rangle|^2}{\langle\psi_1, \psi_1\rangle\,\langle\psi_2, \psi_2\rangle} .
\end{equation}
\begin{itemize}
  \item Scale and phase drop out: replacing $\psi_1$ by $\gamma\psi_1$ multiplies
        numerator and denominator both by $|\gamma|^2$.
  \item By the Cauchy--Schwarz inequality it lies between 0 and 1. It is the squared
        cosine of the angle between the two vectors.
  \item $\mathrm{MAC} = 1$ for the same shape.
  \item $\mathrm{MAC} = 0$ for $M_r$-orthogonal vectors, which by
        Section~\ref{sec:orthogonality} holds for two different bands at the same
        $\mathbf{k}$.
\end{itemize}
Two vectors at neighbouring $\mathbf{k}$ come with two slightly different $M_r$. The
measurements below use $M_r$ of the first of the two points; the difference
vanishes as the step shrinks.

\subsection{How the step size enters}
Normalize $\langle\psi,\psi\rangle = 1$ and write the vector at the next
wavevector as
\begin{equation}
  \psi(\mathbf{k} + \delta\mathbf{k}) \approx \psi + \delta\,(\alpha\,\psi + p),
  \qquad \langle\psi, p\rangle = 0,
\end{equation}
with $\delta$ the step length, $\alpha$ a number and $p$ the part of the change
perpendicular to $\psi$. Then
\begin{align}
  \langle\psi, \psi(\mathbf{k} + \delta\mathbf{k})\rangle &= 1 + \delta\alpha, \\
  \langle\psi(\mathbf{k} + \delta\mathbf{k}), \psi(\mathbf{k} + \delta\mathbf{k})\rangle &= |1 + \delta\alpha|^2 + \delta^2\langle p, p\rangle ,
\end{align}
and so
\begin{equation}
  \mathrm{MAC} = \frac{|1 + \delta\alpha|^2}{|1 + \delta\alpha|^2 + \delta^2\langle p,p\rangle}
  \approx 1 - \delta^2\langle p, p\rangle .
\end{equation}
The part along $\psi$ only changes scale and phase and drops out. So
$1 - \mathrm{MAC} \propto \delta^2$, and halving the step should reduce it by a
factor of 4. Measured for band 1 at its worst step over the whole path (the point
number refers to the segments $\Gamma$--$X$ and $X$--$M$; $M$--$\Gamma$ gets
proportionally more):

\begin{center}
\begin{tabular}{|l|l|l|l|}
\hline
Points/segment & Wavevectors & Lowest MAC, band 1 & $1 - \mathrm{MAC}$ \\
\hline
20 & 69 & 0.9622 & 0.0378 \\
40 & 138 & 0.9898 & 0.0102 \\
80 & 274 & 0.9974 & 0.0026 \\
\hline
\end{tabular}
\end{center}

The ratios are $3.7$ and $3.9$, as expected. At 20 points per segment this worst
step lies just before $X$ on $\Gamma$--$X$, not at the resonator crossing on
$M$--$\Gamma$.

\subsection{Why only neighbours are compared}
Band 1 at $\Gamma$ against band 1 at $M$ has $\mathrm{MAC} = 0.35$, because the mode
legitimately changes from $80\,\%$ plate to $5\,\%$ plate. A low value over a long
distance means nothing. Between neighbouring points on $M$--$\Gamma$ the same band
never drops below $0.9974$ at 40 points per segment, even though it changes its
character completely along that path.

\subsection{Which vector and which weighting}
Lowest neighbouring MAC of band 1 at 40 points per segment:

\begin{center}
\begin{tabular}{|l|l|}
\hline
Variant & Lowest value \\
\hline
Full vector, $M_r$-weighted & 0.990 \\
Plate entries only, $M_r$-weighted & 0.969 \\
Full vector, unweighted & 0.619 \\
\hline
\end{tabular}
\end{center}

The MAC normalizes, so what is left of the plate motion in a resonator-dominated mode
still gives a usable comparison, only slightly worse. The clearly bad variant is
leaving out $M_r$: it counts every entry equally, whether deflection, rotation or
resonator, and falls to $0.62$ on a band that is perfectly smooth. Stage 5 therefore
uses the full vector including $u$, weighted with $M_r$.

% ---------------------------------------------------------------------
\section{How fast a vector turns at an avoided crossing}
\label{sec:mixing_angle}

This section takes the $2\times 2$ matrix of the two-level model apart, to see how
the eigenvector turns through an avoided crossing.

\subsection{The matrix we start from}
For one plate wave and one resonator, in mass-normalized coordinates, the eigenvector
is $(q, v)$ with $q$ the plate wave amplitude and $v = \sqrt{m}\,u$, and
Chapter~\ref{ch:local_resonance} gave
\begin{equation}
  H = \begin{pmatrix} H_{11} & c \\ c & H_{22} \end{pmatrix}
    = \begin{pmatrix} \lambda_p + \mu\lambda_r & -\sqrt{\mu}\,\lambda_r \\ -\sqrt{\mu}\,\lambda_r & \lambda_r \end{pmatrix} .
\end{equation}

\subsection{Mean and half difference}
Define the mean and the half difference of the diagonal,
\begin{equation}
  \bar{H} = \frac{H_{11} + H_{22}}{2}, \qquad \Delta = \frac{H_{11} - H_{22}}{2} .
\end{equation}
Adding and subtracting these gives $H_{11} = \bar{H} + \Delta$ and
$H_{22} = \bar{H} - \Delta$, so
\begin{equation}
  H = \bar{H}\, I + B, \qquad B = \begin{pmatrix} \Delta & c \\ c & -\Delta \end{pmatrix} .
\end{equation}
Inserting the entries:

\begin{center}
\begin{tabular}{|l|l|l|}
\hline
Symbol & In terms of $H$ & For plate wave and resonator \\
\hline
$\bar{H}$ & $(H_{11} + H_{22})/2$ & $\tfrac12\left[\lambda_p + \lambda_r(1+\mu)\right]$ \\
$\Delta$ & $(H_{11} - H_{22})/2$ & $\tfrac12\left[\lambda_p - \lambda_r(1-\mu)\right]$ \\
$c$ & off-diagonal entry & $-\sqrt{\mu}\,\lambda_r$ \\
\hline
\end{tabular}
\end{center}

\paragraph{Why split it.} If $Bz = rz$ for a vector $z$, then
$Hz = \bar{H}z + Bz = (\bar{H} + r)z$. The part $\bar{H}I$ only shifts the
eigenvalues. \emph{The eigenvectors depend on $B$ alone}, so only $\Delta$ and $c$
matter. $\Delta = 0$ at $\lambda_p = \lambda_r(1-\mu)$, where the diagonal entries
are equal: the centre of the avoided crossing.

\subsection{Eigenvalues}
\begin{align}
  \det(B - rI) &= (\Delta - r)(-\Delta - r) - c^2 = r^2 - \Delta^2 - c^2 = 0, \\
  \Rightarrow\quad r &= \pm R, \qquad R = \sqrt{\Delta^2 + c^2} .
\end{align}
The eigenvalues of $H$ are $\bar{H} \pm R$. This is the familiar two-level formula
\begin{equation}
  \lambda = \frac{H_{11} + H_{22}}{2} \pm \sqrt{\left(\frac{H_{11} - H_{22}}{2}\right)^2 + |c|^2} .
\end{equation}
Their sum is $2\bar{H} = \lambda_p + \lambda_r(1+\mu)$ and their product is
\begin{equation}
  \bar{H}^2 - R^2 = \bar{H}^2 - \Delta^2 - c^2 = H_{11}H_{22} - c^2
  = (\lambda_p + \mu\lambda_r)\lambda_r - \mu\lambda_r^2 = \lambda_p\lambda_r .
\end{equation}

\subsection{Eigenvectors and the mixing angle}
$(\Delta, c)$ is a point in a plane at distance $R$ from the origin. Call its polar
angle $2\chi$:
\begin{equation}
  \Delta = R\cos 2\chi, \qquad c = R\sin 2\chi, \qquad \tan 2\chi = \frac{c}{\Delta} .
\end{equation}
This is always possible because $\Delta^2 + c^2 = R^2$. Then
\begin{equation}
  B = R \begin{pmatrix} \cos 2\chi & \sin 2\chi \\ \sin 2\chi & -\cos 2\chi \end{pmatrix} .
\end{equation}
Apply $B$ to $(\cos\chi, \sin\chi)$ and use the addition theorems for
$\cos(2\chi - \chi)$ and $\sin(2\chi - \chi)$:
\begin{align}
  \text{row 1:} \quad & R(\cos 2\chi\cos\chi + \sin 2\chi\sin\chi) = R\cos(2\chi - \chi) = R\cos\chi, \\
  \text{row 2:} \quad & R(\sin 2\chi\cos\chi - \cos 2\chi\sin\chi) = R\sin(2\chi - \chi) = R\sin\chi .
\end{align}
So $(\cos\chi, \sin\chi)$ belongs to $+R$. Applying $B$ to $(-\sin\chi, \cos\chi)$
in the same way:
\begin{align}
  \text{row 1:} \quad & R(-\cos 2\chi\sin\chi + \sin 2\chi\cos\chi) = R\sin\chi = -R\,(-\sin\chi), \\
  \text{row 2:} \quad & R(-\sin 2\chi\sin\chi - \cos 2\chi\cos\chi) = -R\cos\chi ,
\end{align}
so it belongs to $-R$. In summary,
\begin{equation}
  \text{upper band: } (q, v) = (\cos\chi, \sin\chi), \qquad
  \text{lower band: } (q, v) = (-\sin\chi, \cos\chi) .
\end{equation}
The plate share is $|q|^2$. With $\cos^2\chi = (1 + \cos 2\chi)/2$ and
$\cos 2\chi = \Delta/R$,
\begin{equation}
  \boxed{\;\begin{aligned}
  \text{upper band: plate share} &= \tfrac12\left(1 + \frac{\Delta}{R}\right), \\
  \text{lower band: plate share} &= \tfrac12\left(1 - \frac{\Delta}{R}\right)
  \end{aligned}\;}
\end{equation}
Because $c < 0$ here, $\chi$ is negative. That only sets the relative sign of $q$ and
$v$; the shares do not depend on it.

\subsection{One worked point: $\Gamma$}
At $\Gamma$, $\lambda_p = 0$, and with $\mu = 0.25$:
\begin{align}
  H_{11} &= 0.25\,\lambda_r, \quad H_{22} = \lambda_r, \quad c = -0.5\,\lambda_r, \\
  \bar{H} &= 0.625\,\lambda_r, \quad \Delta = -0.375\,\lambda_r, \\
  R &= \sqrt{0.1406 + 0.25}\,\lambda_r = 0.625\,\lambda_r .
\end{align}
The eigenvalues are $\bar{H} \pm R = 0$ and $1.25\,\lambda_r$, the known $0$ and
$\lambda_r(1+\mu)$. With $\Delta/R = -0.6$, the lower band has plate share
$\tfrac12(1 + 0.6) = 80\,\%$ and the upper band $\tfrac12(1 - 0.6) = 20\,\%$.
Further, $2\chi = -126.9^\circ$, so $\chi = -63.4^\circ$.

This agrees with the ratio formula of Chapter~\ref{ch:local_resonance},
$q/v = (\lambda_r - \lambda)/(\sqrt{\mu}\,\lambda_r)$: for the upper band it gives
$-0.25/0.5 = -0.5$, and $\cot(-63.4^\circ) = -0.5$.

It also explains why band 1 starts at $80\,\%$ and not at $100\,\%$. Since $\lambda_p$
cannot go below 0, $\Delta$ cannot go below $-0.375\,\lambda_r$, and the far end of
the turn is never reached.

\subsection{Along $M$--$\Gamma$: model against FE}
Plate shares in percent, for bands 1 and 2 (B1, B2), one-wave model (mod.)\ against
the FE cell:

\begin{center}
\begin{tabular}{|l|l|l|l|l|l|l|l|}
\hline
$k$ & $\lambda_p/\lambda_r$ & $\Delta/\lambda_r$ & $\chi$ & B1 mod. & B1 FE & B2 mod. & B2 FE \\
\hline
0.018 & 0.000 & $-0.375$ & $-63.4^\circ$ & 80.0 & 80.0 & 20.0 & 20.4 \\
0.509 & 0.267 & $-0.242$ & $-57.9^\circ$ & 71.8 & 70.3 & 28.2 & 31.2 \\
0.649 & 0.705 & $-0.023$ & $-46.3^\circ$ & 52.3 & 46.3 & 47.7 & 56.1 \\
0.719 & 1.060 & $0.155$ & $-36.4^\circ$ & 35.2 & 27.7 & 64.8 & 75.2 \\
0.789 & 1.536 & $0.393$ & $-25.9^\circ$ & 19.1 & 14.2 & 80.9 & 89.2 \\
0.930 & 2.943 & $1.096$ & $-12.3^\circ$ & 4.5 & 5.6 & 95.5 & 98.9 \\
1.000 & 3.928 & $1.589$ & $-8.7^\circ$ & 2.3 & 5.1 & 97.7 & 100.0 \\
\hline
\end{tabular}
\end{center}

The model has $\Delta = 0$ at $k \approx 0.667$, and the FE band 1 passes $50\,\%$
shortly before $k = 0.649$. The agreement fades towards $M$, as it should: at $M$
four plate waves take part, and this matrix holds only one. The band frequencies show
the same (band 1 at $M$: $1153\,\mathrm{Hz}$ in the model, $1031\,\mathrm{Hz}$ in
FE). With these values $\chi$ runs from $-63.4^\circ$ at $\Gamma$ to $-8.7^\circ$ at
$M$, about $55^\circ$ in total, not the idealized full turn of $90^\circ$.

\subsection{From the angle to the MAC}
Two normalized vectors at angles $\chi_1$ and $\chi_2$ have the inner product
$\cos\chi_1\cos\chi_2 + \sin\chi_1\sin\chi_2 = \cos(\chi_1 - \chi_2)$, so
\begin{equation}
  \mathrm{MAC} = \cos^2(\chi_1 - \chi_2) .
\end{equation}
To see how fast $\chi$ turns, differentiate $2\chi = \arctan(c/\Delta)$ with respect
to $\Delta$:
\begin{equation}
  2\,\frac{\mathrm{d}\chi}{\mathrm{d}\Delta}
  = \frac{1}{1 + c^2/\Delta^2}\cdot\left(-\frac{c}{\Delta^2}\right)
  = -\frac{c}{\Delta^2 + c^2}
  = -\frac{c}{R^2} .
\end{equation}
The turn is fastest at $\Delta = 0$, where $R = |c|$ and
$|\mathrm{d}\chi/\mathrm{d}\Delta| = 1/(2|c|)$. Since $\Delta$ contains $\lambda_p/2$,
\begin{equation}
  \left|\frac{\mathrm{d}\chi}{\mathrm{d}\lambda_p}\right|_{\Delta = 0}
  = \frac{1}{4|c|} = \frac{1}{4\sqrt{\mu}\,\lambda_r} .
\end{equation}
Half of the full $90^\circ$ turn happens while $|\Delta| < |c|$: at $\Delta = |c|$ one
has $2\chi = -45^\circ$, at $\Delta = -|c|$ one has $2\chi = -135^\circ$, so $\chi$
changes by $45^\circ$. In terms of $\lambda_p$ this window has the width
$4\sqrt{\mu}\,\lambda_r$: $2.0\,\lambda_r$ for $\mu = 0.25$, $0.4\,\lambda_r$ for
$\mu = 0.01$. A lighter resonator has a narrower crossing and needs smaller steps.

Measured for band 1 on $M$--$\Gamma$, with the spacing $0.0175\,\pi/a$ of the 40-point
path (57 steps on this segment):

\begin{center}
\begin{tabular}{|l|l|l|}
\hline
 & Lowest MAC on $M$--$\Gamma$ & at $k$ \\
\hline
Model, $\cos^2(\chi_1 - \chi_2)$ & 0.99778 & 0.737 \\
FE & 0.99736 & 0.702 \\
\hline
\end{tabular}
\end{center}

The minimum lies a little past $\Delta = 0$ because $\lambda_p$ also changes faster
with $k$ as $k$ grows. The change from $69\,\%$ to $24\,\%$ plate takes 12 steps, and
the change from $80\,\%$ to $5\,\%$ takes practically all 57.

\paragraph{The baseline of the FE value.}
Close to $\Gamma$ the FE MAC stays at $0.99961$ while the model gives $1.00000$. The
model uses ``the plate wave at this $\mathbf{k}$'' as a basis vector, so it cannot see
that the pattern $e^{i\mathbf{k}\cdot\mathbf{x}}$ itself changes between neighbouring
points. One step changes the phase at $\mathbf{x}$ by $\theta(\mathbf{x}) =
\delta\mathbf{k}\cdot\mathbf{x}$, and $1 - \mathrm{MAC}$ is approximately the variance
of $\theta$ over the mass of the mode. With $\delta k_x = \delta k_y = 0.0175\,\pi/a$,
a variance of $a^2/12$ of each coordinate over the cell, and $80\,\%$ of the mass
motion in the plate (the resonator sits at the mean position):
\begin{equation}
  1 - \mathrm{MAC} \approx 0.8\cdot (0.0175\,\pi)^2\cdot\frac{2}{12} \approx 0.0004 ,
\end{equation}
which matches the measured $1 - 0.99961 = 0.00039$.

\subsection{The other two-level matrix: two plate waves and a rigid mass}
Chapter~\ref{ch:two_level} used the same structure with
$H_{11} = \lambda_1^0 - \mu\lambda$, $H_{22} = \lambda_2^0 - \mu\lambda$ and
$c = -\mu\lambda$. Then
\begin{equation}
  \bar{H} = \frac{\lambda_1^0 + \lambda_2^0}{2} - \mu\lambda, \qquad
  \Delta = \frac{\lambda_1^0 - \lambda_2^0}{2},
\end{equation}
because the mass loading cancels in the difference. At $X$, $\Delta = 0$, so
$2\chi = -90^\circ$ and $\chi = -45^\circ$. That gives the eigenvectors
$(1, -1)/\sqrt{2}$ and $(1, 1)/\sqrt{2}$ of that chapter, with $a_H = \bar{H}$ and
$b_H = c$.

% ---------------------------------------------------------------------
\section{The three situations a tracker meets}

Lowest neighbouring MAC between \emph{sorted} bands along the whole path:

\begin{center}
\begin{tabular}{|l|l|l|l|p{0.32\textwidth}|}
\hline
Band & 20 points & 40 points & 80 points & Where \\
\hline
1 & 0.9622 & 0.9898 & 0.9974 & smooth band \\
2 & 0.2301 & 0.2325 & 0.2322 & next to $M$, three modes at $2379\,\mathrm{Hz}$ \\
3 & 0.1755 & 0.1954 & 0.2117 & next to $\Gamma$, three modes at $4730\,\mathrm{Hz}$ \\
4 & 0.0601 & 0.0610 & 0.0615 & next to $X$, three modes at $5881\,\mathrm{Hz}$ \\
5 & 0.0000 & 0.0000 & 0.0000 & near $k \approx 0.33$ on $M$--$\Gamma$, about $6.5\,\mathrm{kHz}$ \\
\hline
\end{tabular}
\end{center}

\paragraph{1. Smooth continuation.} Band 1: the MAC is close to 1 and improves with a
finer step.

\paragraph{2. A degenerate point.} Bands 2 to 4: the value barely changes with a finer
step. That is the signature of a degeneracy and not of a coarse step. At the
degenerate point the solver picks an arbitrary basis of the group
(Section~\ref{sec:orthogonality}), while just beside it the modes are distinct. Near
$M$, MAC of the three neighbouring modes against the three modes at $M$ (40 points per
segment):

\begin{center}
\begin{tabular}{|l|l|l|l|l|}
\hline
Before $M$ & $M$, mode 1 & $M$, mode 2 & $M$, mode 3 & Row sum \\
\hline
$2321\,\mathrm{Hz}$ & 0.497 & 0.045 & 0.457 & 1.000 \\
$2372\,\mathrm{Hz}$ & 0.017 & 0.882 & 0.072 & 0.971 \\
$2439\,\mathrm{Hz}$ & 0.485 & 0.047 & 0.468 & 1.000 \\
\hline
\end{tabular}
\end{center}

No single entry tells the continuation, but the row sums are close to 1. The reason:
let $e_1, \dots, e_g$ be an $M_r$-orthonormal basis of the group and $\psi$ a vector
with $\langle\psi,\psi\rangle = 1$. Then
\begin{equation}
  \sum_{j=1}^{g} \mathrm{MAC}(\psi, e_j) = \sum_{j=1}^{g} |\langle e_j, \psi\rangle|^2,
\end{equation}
which by Pythagoras is the squared length of the part of $\psi$ that lies inside the
group. It equals 1 exactly if $\psi$ lies entirely in the group, and it does not depend
on which basis the solver chose. The neighbouring mode sits at a slightly different
$\mathbf{k}$, so the sum is only approximately 1. Degenerate groups therefore have to
be compared as a whole, and all members of the group must have been computed.

\paragraph{3. A true crossing.} Two sorted bands swap identities: the sorted MAC is
about 0 and the MAC with the other band is about 1. In the sweep, bands 5 and 6 near
$k \approx 0.33$ on $M$--$\Gamma$ (about $6.5\,\mathrm{kHz}$, above the $5\,\mathrm{kHz}$
range) show a MAC of $0.0000$. That is consistent with two modes of opposite symmetry
about the diagonal, which have $c = 0$; the symmetry itself was not checked.

% ---------------------------------------------------------------------
\section{Tracking and gap finding are different tasks}

\emph{Tracking} follows a mode's identity from one $\mathbf{k}$ to the next, using the
MAC between neighbours. It answers ``which curve is this dot on?''

\emph{Gap finding} asks whether any dot lies in a frequency interval. It needs only the
frequencies sorted at each $\mathbf{k}$: the gap between sorted bands $n$ and $n+1$
runs from the maximum of band $n$ to the minimum of band $n+1$.

A tracked band may cross another one, so the $n$-th tracked band is not the $n$-th
sorted band. Suppose tracked curves $A$ and $B$ cross and a third band $C$ lies above
both. The sorted band below $C$ is the upper envelope of $A$ and $B$ at each
$\mathbf{k}$, so the correct lower gap edge is the maximum of that envelope. Taking the
curve $B$ alone when $A$ rises higher later on the path gives a value that is too low:
the reported gap starts too low and contains dots of $A$. Stage 5 therefore keeps plain
sorting available.

% ---------------------------------------------------------------------
\section{Check questions with answers}

\paragraph{Question 1.} \emph{Band 3 near $\Gamma$ has a neighbouring MAC of about 0.2,
and it stays there when the number of $\mathbf{k}$-points is doubled. What do you
conclude, and what would you do differently for that step?}

\emph{Answer.} Band 3 meets two other bands at $\Gamma$, and the solver returns an
arbitrary basis of that group. The MAC has to be summed over the whole group, which
then gives approximately 1. A degeneracy has to be told apart from a true crossing
between two sampled points: there the continuation is well defined, the MAC with the
partner band is about 1 and the sorted labels simply swap. Here the MAC is about 0.2,
neither 0 nor 1, and refining the step does not change it, which points to a degenerate
point where individual vectors mean nothing. The sum is only approximately 1 because the
neighbouring vector sits at a slightly different $\mathbf{k}$ (measured: $0.97$ to
$1.00$), and it needs every member of the group, which is why the solver uses guard
eigenpairs. For that step: match the group as a whole instead of one vector to one
vector. At the degenerate dot all members have the same frequency, so the plotted
curve does not depend on which label goes where.

\paragraph{Question 2.} \emph{Band 1 changes from $80\,\%$ plate to $5\,\%$ plate along
$M$--$\Gamma$, yet the neighbouring MAC stays close to 1 there. Why is this not a
contradiction, and which change in the unit cell would make this crossing harder to
track?}

\emph{Answer.} The change is gradual, spread over essentially all 57 steps of the
segment ($69\,\%$ to $24\,\%$ alone takes 12 steps), and each step only compares two
nearly identical shapes. A lighter resonator makes the crossing narrower, because the
coupling is $c = -\sqrt{\mu}\,\lambda_r$; the turn then happens within fewer steps, and
at the same step size the MAC dips further from 1.

\paragraph{Question 3.} \emph{Two tracked bands overlap in frequency, with a true crossing
between them. Could there still be a complete gap between sorted bands $n$ and $n+1$ at
that crossing frequency, and which data decides it?}

\emph{Answer.} No. At a true crossing the sorted bands $n$ and $n+1$ share a dot at the
crossing frequency $f_c$, so $\max(\text{band } n) \ge f_c \ge \min(\text{band } n+1)$
and there is no gap between that pair. The data to decide this is the sorted frequencies
at every $\mathbf{k}$; no eigenvectors are needed. The point of the question is the
reverse danger: using tracked bands for gap edges can report a gap that contains dots of
another band, as shown in the previous section.

\paragraph{Follow-up: how does the matrix $H$ of the two-level model translate into the
notation with mean, half difference and mixing angle?}
\emph{Answer.} This is worked out in Section~\ref{sec:mixing_angle}: $H_{11}$ and
$H_{22}$ are replaced by $\bar{H} \pm \Delta$, the off-diagonal $c$ stays the same, and
the eigenvectors depend only on $c/\Delta$.

% =====================================================================

\chapter{Eigenvectors, resonator participation and the MAC}
\label{ch:eigenvectors}

So far only the eigenvalues of the Bloch problem were used. This chapter looks at
the eigenvectors: what they contain, how much of a mode sits in the resonator, how
two modes at neighbouring wavevectors are compared with the Modal Assurance
Criterion (MAC), and which situations a mode tracker has to handle. All FE numbers
are measured on the production cell: $a = 0.1\,\mathrm{m}$, $h = 5\,\mathrm{mm}$,
steel, $10\times 10$ Lagrange degree 2 elements, one resonator at the cell centre
with $f_r = 1200\,\mathrm{Hz}$ and $\mu = 0.25$.

\begin{center}
\begin{tabular}{|l|p{0.68\textwidth}|}
\hline
Symbol & Meaning \\
\hline
$K_r$, $M_r$ & reduced stiffness and mass matrices at one $\mathbf{k}$, $K_r = T^H K T$, $M_r = T^H M T$ \\
$T(\mathbf{k})$ & maps independent DOFs to all DOFs, with phase factors \\
$\lambda = \omega^2$ & eigenvalue of one dot in the diagram \\
$\psi$ & reduced eigenvector of one dot, the column returned by the solver \\
$u$ & last entry of $\psi$: the resonator displacement \\
$m$, $k_s$ & resonator mass and spring stiffness \\
$\lambda_r = k_s/m$ & the resonator's own eigenvalue \\
$\lambda_p$ & bare plate eigenvalue of one plate wave \\
$\mu = m/(\rho h A)$ & resonator mass over plate mass per cell \\
$\langle \psi_1, \psi_2\rangle$ & mass-weighted inner product $\psi_1^H M_r \psi_2$ \\
$H$ & $2\times 2$ standard-form matrix of the two-level model \\
$q$, $v$ & entries of the eigenvector of $H$: plate wave amplitude and $v = \sqrt{m}\,u$ \\
$\bar{H}$, $\Delta$ & mean and half difference of the diagonal entries of $H$ \\
$c$ & off-diagonal entry of $H$ \\
$R$ & $\sqrt{\Delta^2 + c^2}$ \\
$\chi$ & mixing angle of the $2\times 2$ eigenvector \\
$k$ (in tables) & position on the diagonal, $k_x = k_y = k\,\pi/a$ \\
\hline
\end{tabular}
\end{center}

% ---------------------------------------------------------------------
\section{What an eigenvector is at one wavevector}

At each $\mathbf{k}$ the solver returns pairs $(\lambda, \psi)$ with
\begin{equation}
  K_r \psi = \lambda M_r \psi .
\end{equation}
Each dot in the diagram owns one $\psi$. Its entries are the plate unknowns
(deflection and two rotations) at the independent nodes, and its last entry is the
resonator displacement $u$. The resonator has no periodic partner, so $T$ acts on it
as the identity, and the last row of $M_r$ contains only $m$ on the diagonal.

$\psi$ is defined only up to a factor: if $\psi$ solves the equation, so does
$\gamma\psi$ for any complex number $\gamma$. Its scale and its overall complex
phase therefore carry no physical meaning, and everything computed from $\psi$ has
to be independent of both.

% ---------------------------------------------------------------------
\section{Orthogonality and what a degeneracy does}
\label{sec:orthogonality}

Take two solutions at the same $\mathbf{k}$:
\begin{equation}
  K_r \psi_1 = \lambda_1 M_r \psi_1, \qquad K_r \psi_2 = \lambda_2 M_r \psi_2 .
\end{equation}
Multiply the first from the left by $\psi_2^H$:
\begin{equation}
  \psi_2^H K_r \psi_1 = \lambda_1\, \psi_2^H M_r \psi_1 .
\end{equation}
Take the conjugate transpose of the second. $K_r$ and $M_r$ are Hermitian and
$\lambda_2$ is real, so $\psi_2^H K_r = \lambda_2\, \psi_2^H M_r$. Multiply from the
right by $\psi_1$:
\begin{equation}
  \psi_2^H K_r \psi_1 = \lambda_2\, \psi_2^H M_r \psi_1 .
\end{equation}
The left sides are equal. Subtracting gives
\begin{equation}
  (\lambda_1 - \lambda_2)\, \psi_2^H M_r \psi_1 = 0 .
\end{equation}
Different eigenvalues therefore mean $\psi_2^H M_r \psi_1 = 0$: the vectors are
orthogonal when weighted with $M_r$. That is why the mass matrix appears in
everything below.

If $\lambda_1 = \lambda_2 = \lambda$, the equation says nothing. Then every
combination is also a solution,
\begin{equation}
  K_r(\alpha_1\psi_1 + \alpha_2\psi_2) = \lambda M_r(\alpha_1\psi_1 + \alpha_2\psi_2),
\end{equation}
so at a degenerate point the solver returns \emph{some arbitrary basis} of the
group. The individual vectors of the group have no meaning of their own; only the
space they span does.

In this cell, degenerate groups appear where several bare plate waves share one
frequency. One combination couples to the resonator and leaves the group; the rest
do not see it. For example, $M$ has four plate waves at $2378\,\mathrm{Hz}$: one
coupled combination moves away and three modes remain at $2379\,\mathrm{Hz}$.

% ---------------------------------------------------------------------
\section{Plate share in the FE model}

The kinetic energy of a mode is $\tfrac12\omega^2\,\psi^H M_r\psi$, and the part in
the resonator is $\tfrac12\omega^2\, m|u|^2$. Hence
\begin{equation}
  \text{plate share} = 1 - \frac{m|u|^2}{\psi^H M_r \psi} .
\end{equation}
Scale and phase of $\psi$ cancel, as required.

\begin{center}
\begin{tabular}{|l|p{0.72\textwidth}|}
\hline
Point & Frequency in Hz: plate share (FE) \\
\hline
$\Gamma$ & $0$: $80\,\%$;\quad $1274$: $20\,\%$;\quad $4730$: $100\,\%$ (three modes);\quad $4883$: $100\,\%$ \\
$X$ & $836$: $31\,\%$;\quad $1195$: $100\,\%$;\quad $1653$: $69\,\%$;\quad $5881$: $100\,\%$ (three modes) \\
$M$ & $1031$: $5\,\%$;\quad $2379$: $100\,\%$ (three modes);\quad $2712$: $95\,\%$;\quad $11567$: $100\,\%$ \\
\hline
\end{tabular}
\end{center}

The simple model of Chapter~\ref{ch:local_resonance} gave $34\,\%$ and $66\,\%$ at
$X$; the shares at $\Gamma$ and $M$ agree with it.

\paragraph{The sum rule for the resonator share.}
Collect the eigenvectors of all bands at one $\mathbf{k}$ as columns of a matrix
$\Psi$, normalized so that $\Psi^H M_r \Psi = I$. Then $\Psi^{-1} = \Psi^H M_r$, so
\begin{equation}
  \Psi\,\Psi^H M_r = I \quad\Rightarrow\quad \Psi\,\Psi^H = M_r^{-1} .
\end{equation}
The resonator has no mass coupling to the plate, so the last diagonal entry of
$M_r^{-1}$ is $1/m$. The last diagonal entry of $\Psi\Psi^H$ is $\sum_n |u_n|^2$,
with $n$ running over all bands. Therefore
\begin{equation}
  \sum_n m|u_n|^2 = 1 .
\end{equation}
With this normalization $m|u_n|^2$ is the resonator share of band $n$, so
\emph{the resonator shares of all bands at one $\mathbf{k}$ add up to exactly
$100\,\%$}. The FE values at the three symmetry points show it: $20 + 80$ at
$\Gamma$, $69 + 31$ at $X$, $95 + 5$ at $M$. The plate shares of bands 1 and 2 alone
need not add up to $100\,\%$: along $M$--$\Gamma$ their sum drifts from $100\,\%$ to
about $104\,\%$ close to $M$, because other bands take part there.

% ---------------------------------------------------------------------
\section{The Modal Assurance Criterion}

\subsection{Definition and properties}
We need a number that says how alike two eigenvectors are and that ignores scale
and phase. With $\langle \psi_1, \psi_2\rangle = \psi_1^H M_r \psi_2$,
\begin{equation}
  \mathrm{MAC}(\psi_1, \psi_2) = \frac{|\langle\psi_1, \psi_2\rangle|^2}{\langle\psi_1, \psi_1\rangle\,\langle\psi_2, \psi_2\rangle} .
\end{equation}
\begin{itemize}
  \item Scale and phase drop out: replacing $\psi_1$ by $\gamma\psi_1$ multiplies
        numerator and denominator both by $|\gamma|^2$.
  \item By the Cauchy--Schwarz inequality it lies between 0 and 1. It is the squared
        cosine of the angle between the two vectors.
  \item $\mathrm{MAC} = 1$ for the same shape.
  \item $\mathrm{MAC} = 0$ for $M_r$-orthogonal vectors, which by
        Section~\ref{sec:orthogonality} holds for two different bands at the same
        $\mathbf{k}$.
\end{itemize}
Two vectors at neighbouring $\mathbf{k}$ come with two slightly different $M_r$. The
measurements below use $M_r$ of the first of the two points; the difference
vanishes as the step shrinks.

\subsection{How the step size enters}
Normalize $\langle\psi,\psi\rangle = 1$ and write the vector at the next
wavevector as
\begin{equation}
  \psi(\mathbf{k} + \delta\mathbf{k}) \approx \psi + \delta\,(\alpha\,\psi + p),
  \qquad \langle\psi, p\rangle = 0,
\end{equation}
with $\delta$ the step length, $\alpha$ a number and $p$ the part of the change
perpendicular to $\psi$. Then
\begin{align}
  \langle\psi, \psi(\mathbf{k} + \delta\mathbf{k})\rangle &= 1 + \delta\alpha, \\
  \langle\psi(\mathbf{k} + \delta\mathbf{k}), \psi(\mathbf{k} + \delta\mathbf{k})\rangle &= |1 + \delta\alpha|^2 + \delta^2\langle p, p\rangle ,
\end{align}
and so
\begin{equation}
  \mathrm{MAC} = \frac{|1 + \delta\alpha|^2}{|1 + \delta\alpha|^2 + \delta^2\langle p,p\rangle}
  \approx 1 - \delta^2\langle p, p\rangle .
\end{equation}
The part along $\psi$ only changes scale and phase and drops out. So
$1 - \mathrm{MAC} \propto \delta^2$, and halving the step should reduce it by a
factor of 4. Measured for band 1 at its worst step over the whole path (the point
number refers to the segments $\Gamma$--$X$ and $X$--$M$; $M$--$\Gamma$ gets
proportionally more):

\begin{center}
\begin{tabular}{|l|l|l|l|}
\hline
Points/segment & Wavevectors & Lowest MAC, band 1 & $1 - \mathrm{MAC}$ \\
\hline
20 & 69 & 0.9622 & 0.0378 \\
40 & 138 & 0.9898 & 0.0102 \\
80 & 274 & 0.9974 & 0.0026 \\
\hline
\end{tabular}
\end{center}

The ratios are $3.7$ and $3.9$, as expected. At 20 points per segment this worst
step lies just before $X$ on $\Gamma$--$X$, not at the resonator crossing on
$M$--$\Gamma$.

\subsection{Why only neighbours are compared}
Band 1 at $\Gamma$ against band 1 at $M$ has $\mathrm{MAC} = 0.35$, because the mode
legitimately changes from $80\,\%$ plate to $5\,\%$ plate. A low value over a long
distance means nothing. Between neighbouring points on $M$--$\Gamma$ the same band
never drops below $0.9974$ at 40 points per segment, even though it changes its
character completely along that path.

\subsection{Which vector and which weighting}
Lowest neighbouring MAC of band 1 at 40 points per segment:

\begin{center}
\begin{tabular}{|l|l|}
\hline
Variant & Lowest value \\
\hline
Full vector, $M_r$-weighted & 0.990 \\
Plate entries only, $M_r$-weighted & 0.969 \\
Full vector, unweighted & 0.619 \\
\hline
\end{tabular}
\end{center}

The MAC normalizes, so what is left of the plate motion in a resonator-dominated mode
still gives a usable comparison, only slightly worse. The clearly bad variant is
leaving out $M_r$: it counts every entry equally, whether deflection, rotation or
resonator, and falls to $0.62$ on a band that is perfectly smooth. Stage 5 therefore
uses the full vector including $u$, weighted with $M_r$.

% ---------------------------------------------------------------------
\section{How fast a vector turns at an avoided crossing}
\label{sec:mixing_angle}

This section takes the $2\times 2$ matrix of the two-level model apart, to see how
the eigenvector turns through an avoided crossing.

\subsection{The matrix we start from}
For one plate wave and one resonator, in mass-normalized coordinates, the eigenvector
is $(q, v)$ with $q$ the plate wave amplitude and $v = \sqrt{m}\,u$, and
Chapter~\ref{ch:local_resonance} gave
\begin{equation}
  H = \begin{pmatrix} H_{11} & c \\ c & H_{22} \end{pmatrix}
    = \begin{pmatrix} \lambda_p + \mu\lambda_r & -\sqrt{\mu}\,\lambda_r \\ -\sqrt{\mu}\,\lambda_r & \lambda_r \end{pmatrix} .
\end{equation}

\subsection{Mean and half difference}
Define the mean and the half difference of the diagonal,
\begin{equation}
  \bar{H} = \frac{H_{11} + H_{22}}{2}, \qquad \Delta = \frac{H_{11} - H_{22}}{2} .
\end{equation}
Adding and subtracting these gives $H_{11} = \bar{H} + \Delta$ and
$H_{22} = \bar{H} - \Delta$, so
\begin{equation}
  H = \bar{H}\, I + B, \qquad B = \begin{pmatrix} \Delta & c \\ c & -\Delta \end{pmatrix} .
\end{equation}
Inserting the entries:

\begin{center}
\begin{tabular}{|l|l|l|}
\hline
Symbol & In terms of $H$ & For plate wave and resonator \\
\hline
$\bar{H}$ & $(H_{11} + H_{22})/2$ & $\tfrac12\left[\lambda_p + \lambda_r(1+\mu)\right]$ \\
$\Delta$ & $(H_{11} - H_{22})/2$ & $\tfrac12\left[\lambda_p - \lambda_r(1-\mu)\right]$ \\
$c$ & off-diagonal entry & $-\sqrt{\mu}\,\lambda_r$ \\
\hline
\end{tabular}
\end{center}

\paragraph{Why split it.} If $Bz = rz$ for a vector $z$, then
$Hz = \bar{H}z + Bz = (\bar{H} + r)z$. The part $\bar{H}I$ only shifts the
eigenvalues. \emph{The eigenvectors depend on $B$ alone}, so only $\Delta$ and $c$
matter. $\Delta = 0$ at $\lambda_p = \lambda_r(1-\mu)$, where the diagonal entries
are equal: the centre of the avoided crossing.

\subsection{Eigenvalues}
\begin{align}
  \det(B - rI) &= (\Delta - r)(-\Delta - r) - c^2 = r^2 - \Delta^2 - c^2 = 0, \\
  \Rightarrow\quad r &= \pm R, \qquad R = \sqrt{\Delta^2 + c^2} .
\end{align}
The eigenvalues of $H$ are $\bar{H} \pm R$. This is the familiar two-level formula
\begin{equation}
  \lambda = \frac{H_{11} + H_{22}}{2} \pm \sqrt{\left(\frac{H_{11} - H_{22}}{2}\right)^2 + |c|^2} .
\end{equation}
Their sum is $2\bar{H} = \lambda_p + \lambda_r(1+\mu)$ and their product is
\begin{equation}
  \bar{H}^2 - R^2 = \bar{H}^2 - \Delta^2 - c^2 = H_{11}H_{22} - c^2
  = (\lambda_p + \mu\lambda_r)\lambda_r - \mu\lambda_r^2 = \lambda_p\lambda_r .
\end{equation}

\subsection{Eigenvectors and the mixing angle}
$(\Delta, c)$ is a point in a plane at distance $R$ from the origin. Call its polar
angle $2\chi$:
\begin{equation}
  \Delta = R\cos 2\chi, \qquad c = R\sin 2\chi, \qquad \tan 2\chi = \frac{c}{\Delta} .
\end{equation}
This is always possible because $\Delta^2 + c^2 = R^2$. Then
\begin{equation}
  B = R \begin{pmatrix} \cos 2\chi & \sin 2\chi \\ \sin 2\chi & -\cos 2\chi \end{pmatrix} .
\end{equation}
Apply $B$ to $(\cos\chi, \sin\chi)$ and use the addition theorems for
$\cos(2\chi - \chi)$ and $\sin(2\chi - \chi)$:
\begin{align}
  \text{row 1:} \quad & R(\cos 2\chi\cos\chi + \sin 2\chi\sin\chi) = R\cos(2\chi - \chi) = R\cos\chi, \\
  \text{row 2:} \quad & R(\sin 2\chi\cos\chi - \cos 2\chi\sin\chi) = R\sin(2\chi - \chi) = R\sin\chi .
\end{align}
So $(\cos\chi, \sin\chi)$ belongs to $+R$. Applying $B$ to $(-\sin\chi, \cos\chi)$
in the same way:
\begin{align}
  \text{row 1:} \quad & R(-\cos 2\chi\sin\chi + \sin 2\chi\cos\chi) = R\sin\chi = -R\,(-\sin\chi), \\
  \text{row 2:} \quad & R(-\sin 2\chi\sin\chi - \cos 2\chi\cos\chi) = -R\cos\chi ,
\end{align}
so it belongs to $-R$. In summary,
\begin{equation}
  \text{upper band: } (q, v) = (\cos\chi, \sin\chi), \qquad
  \text{lower band: } (q, v) = (-\sin\chi, \cos\chi) .
\end{equation}
The plate share is $|q|^2$. With $\cos^2\chi = (1 + \cos 2\chi)/2$ and
$\cos 2\chi = \Delta/R$,
\begin{equation}
  \boxed{\;\begin{aligned}
  \text{upper band: plate share} &= \tfrac12\left(1 + \frac{\Delta}{R}\right), \\
  \text{lower band: plate share} &= \tfrac12\left(1 - \frac{\Delta}{R}\right)
  \end{aligned}\;}
\end{equation}
Because $c < 0$ here, $\chi$ is negative. That only sets the relative sign of $q$ and
$v$; the shares do not depend on it.

\subsection{One worked point: $\Gamma$}
At $\Gamma$, $\lambda_p = 0$, and with $\mu = 0.25$:
\begin{align}
  H_{11} &= 0.25\,\lambda_r, \quad H_{22} = \lambda_r, \quad c = -0.5\,\lambda_r, \\
  \bar{H} &= 0.625\,\lambda_r, \quad \Delta = -0.375\,\lambda_r, \\
  R &= \sqrt{0.1406 + 0.25}\,\lambda_r = 0.625\,\lambda_r .
\end{align}
The eigenvalues are $\bar{H} \pm R = 0$ and $1.25\,\lambda_r$, the known $0$ and
$\lambda_r(1+\mu)$. With $\Delta/R = -0.6$, the lower band has plate share
$\tfrac12(1 + 0.6) = 80\,\%$ and the upper band $\tfrac12(1 - 0.6) = 20\,\%$.
Further, $2\chi = -126.9^\circ$, so $\chi = -63.4^\circ$.

This agrees with the ratio formula of Chapter~\ref{ch:local_resonance},
$q/v = (\lambda_r - \lambda)/(\sqrt{\mu}\,\lambda_r)$: for the upper band it gives
$-0.25/0.5 = -0.5$, and $\cot(-63.4^\circ) = -0.5$.

It also explains why band 1 starts at $80\,\%$ and not at $100\,\%$. Since $\lambda_p$
cannot go below 0, $\Delta$ cannot go below $-0.375\,\lambda_r$, and the far end of
the turn is never reached.

\subsection{Along $M$--$\Gamma$: model against FE}
Plate shares in percent, for bands 1 and 2 (B1, B2), one-wave model (mod.)\ against
the FE cell:

\begin{center}
\begin{tabular}{|l|l|l|l|l|l|l|l|}
\hline
$k$ & $\lambda_p/\lambda_r$ & $\Delta/\lambda_r$ & $\chi$ & B1 mod. & B1 FE & B2 mod. & B2 FE \\
\hline
0.018 & 0.000 & $-0.375$ & $-63.4^\circ$ & 80.0 & 80.0 & 20.0 & 20.4 \\
0.509 & 0.267 & $-0.242$ & $-57.9^\circ$ & 71.8 & 70.3 & 28.2 & 31.2 \\
0.649 & 0.705 & $-0.023$ & $-46.3^\circ$ & 52.3 & 46.3 & 47.7 & 56.1 \\
0.719 & 1.060 & $0.155$ & $-36.4^\circ$ & 35.2 & 27.7 & 64.8 & 75.2 \\
0.789 & 1.536 & $0.393$ & $-25.9^\circ$ & 19.1 & 14.2 & 80.9 & 89.2 \\
0.930 & 2.943 & $1.096$ & $-12.3^\circ$ & 4.5 & 5.6 & 95.5 & 98.9 \\
1.000 & 3.928 & $1.589$ & $-8.7^\circ$ & 2.3 & 5.1 & 97.7 & 100.0 \\
\hline
\end{tabular}
\end{center}

The model has $\Delta = 0$ at $k \approx 0.667$, and the FE band 1 passes $50\,\%$
shortly before $k = 0.649$. The agreement fades towards $M$, as it should: at $M$
four plate waves take part, and this matrix holds only one. The band frequencies show
the same (band 1 at $M$: $1153\,\mathrm{Hz}$ in the model, $1031\,\mathrm{Hz}$ in
FE). With these values $\chi$ runs from $-63.4^\circ$ at $\Gamma$ to $-8.7^\circ$ at
$M$, about $55^\circ$ in total, not the idealized full turn of $90^\circ$.

\subsection{From the angle to the MAC}
Two normalized vectors at angles $\chi_1$ and $\chi_2$ have the inner product
$\cos\chi_1\cos\chi_2 + \sin\chi_1\sin\chi_2 = \cos(\chi_1 - \chi_2)$, so
\begin{equation}
  \mathrm{MAC} = \cos^2(\chi_1 - \chi_2) .
\end{equation}
To see how fast $\chi$ turns, differentiate $2\chi = \arctan(c/\Delta)$ with respect
to $\Delta$:
\begin{equation}
  2\,\frac{\mathrm{d}\chi}{\mathrm{d}\Delta}
  = \frac{1}{1 + c^2/\Delta^2}\cdot\left(-\frac{c}{\Delta^2}\right)
  = -\frac{c}{\Delta^2 + c^2}
  = -\frac{c}{R^2} .
\end{equation}
The turn is fastest at $\Delta = 0$, where $R = |c|$ and
$|\mathrm{d}\chi/\mathrm{d}\Delta| = 1/(2|c|)$. Since $\Delta$ contains $\lambda_p/2$,
\begin{equation}
  \left|\frac{\mathrm{d}\chi}{\mathrm{d}\lambda_p}\right|_{\Delta = 0}
  = \frac{1}{4|c|} = \frac{1}{4\sqrt{\mu}\,\lambda_r} .
\end{equation}
Half of the full $90^\circ$ turn happens while $|\Delta| < |c|$: at $\Delta = |c|$ one
has $2\chi = -45^\circ$, at $\Delta = -|c|$ one has $2\chi = -135^\circ$, so $\chi$
changes by $45^\circ$. In terms of $\lambda_p$ this window has the width
$4\sqrt{\mu}\,\lambda_r$: $2.0\,\lambda_r$ for $\mu = 0.25$, $0.4\,\lambda_r$ for
$\mu = 0.01$. A lighter resonator has a narrower crossing and needs smaller steps.

Measured for band 1 on $M$--$\Gamma$, with the spacing $0.0175\,\pi/a$ of the 40-point
path (57 steps on this segment):

\begin{center}
\begin{tabular}{|l|l|l|}
\hline
 & Lowest MAC on $M$--$\Gamma$ & at $k$ \\
\hline
Model, $\cos^2(\chi_1 - \chi_2)$ & 0.99778 & 0.737 \\
FE & 0.99736 & 0.702 \\
\hline
\end{tabular}
\end{center}

The minimum lies a little past $\Delta = 0$ because $\lambda_p$ also changes faster
with $k$ as $k$ grows. The change from $69\,\%$ to $24\,\%$ plate takes 12 steps, and
the change from $80\,\%$ to $5\,\%$ takes practically all 57.

\paragraph{The baseline of the FE value.}
Close to $\Gamma$ the FE MAC stays at $0.99961$ while the model gives $1.00000$. The
model uses ``the plate wave at this $\mathbf{k}$'' as a basis vector, so it cannot see
that the pattern $e^{i\mathbf{k}\cdot\mathbf{x}}$ itself changes between neighbouring
points. One step changes the phase at $\mathbf{x}$ by $\theta(\mathbf{x}) =
\delta\mathbf{k}\cdot\mathbf{x}$, and $1 - \mathrm{MAC}$ is approximately the variance
of $\theta$ over the mass of the mode. With $\delta k_x = \delta k_y = 0.0175\,\pi/a$,
a variance of $a^2/12$ of each coordinate over the cell, and $80\,\%$ of the mass
motion in the plate (the resonator sits at the mean position):
\begin{equation}
  1 - \mathrm{MAC} \approx 0.8\cdot (0.0175\,\pi)^2\cdot\frac{2}{12} \approx 0.0004 ,
\end{equation}
which matches the measured $1 - 0.99961 = 0.00039$.

\subsection{The other two-level matrix: two plate waves and a rigid mass}
Chapter~\ref{ch:two_level} used the same structure with
$H_{11} = \lambda_1^0 - \mu\lambda$, $H_{22} = \lambda_2^0 - \mu\lambda$ and
$c = -\mu\lambda$. Then
\begin{equation}
  \bar{H} = \frac{\lambda_1^0 + \lambda_2^0}{2} - \mu\lambda, \qquad
  \Delta = \frac{\lambda_1^0 - \lambda_2^0}{2},
\end{equation}
because the mass loading cancels in the difference. At $X$, $\Delta = 0$, so
$2\chi = -90^\circ$ and $\chi = -45^\circ$. That gives the eigenvectors
$(1, -1)/\sqrt{2}$ and $(1, 1)/\sqrt{2}$ of that chapter, with $a_H = \bar{H}$ and
$b_H = c$.

% ---------------------------------------------------------------------
\section{The three situations a tracker meets}

Lowest neighbouring MAC between \emph{sorted} bands along the whole path:

\begin{center}
\begin{tabular}{|l|l|l|l|p{0.32\textwidth}|}
\hline
Band & 20 points & 40 points & 80 points & Where \\
\hline
1 & 0.9622 & 0.9898 & 0.9974 & smooth band \\
2 & 0.2301 & 0.2325 & 0.2322 & next to $M$, three modes at $2379\,\mathrm{Hz}$ \\
3 & 0.1755 & 0.1954 & 0.2117 & next to $\Gamma$, three modes at $4730\,\mathrm{Hz}$ \\
4 & 0.0601 & 0.0610 & 0.0615 & next to $X$, three modes at $5881\,\mathrm{Hz}$ \\
5 & 0.0000 & 0.0000 & 0.0000 & near $k \approx 0.33$ on $M$--$\Gamma$, about $6.5\,\mathrm{kHz}$ \\
\hline
\end{tabular}
\end{center}

\paragraph{1. Smooth continuation.} Band 1: the MAC is close to 1 and improves with a
finer step.

\paragraph{2. A degenerate point.} Bands 2 to 4: the value barely changes with a finer
step. That is the signature of a degeneracy and not of a coarse step. At the
degenerate point the solver picks an arbitrary basis of the group
(Section~\ref{sec:orthogonality}), while just beside it the modes are distinct. Near
$M$, MAC of the three neighbouring modes against the three modes at $M$ (40 points per
segment):

\begin{center}
\begin{tabular}{|l|l|l|l|l|}
\hline
Before $M$ & $M$, mode 1 & $M$, mode 2 & $M$, mode 3 & Row sum \\
\hline
$2321\,\mathrm{Hz}$ & 0.497 & 0.045 & 0.457 & 1.000 \\
$2372\,\mathrm{Hz}$ & 0.017 & 0.882 & 0.072 & 0.971 \\
$2439\,\mathrm{Hz}$ & 0.485 & 0.047 & 0.468 & 1.000 \\
\hline
\end{tabular}
\end{center}

No single entry tells the continuation, but the row sums are close to 1. The reason:
let $e_1, \dots, e_g$ be an $M_r$-orthonormal basis of the group and $\psi$ a vector
with $\langle\psi,\psi\rangle = 1$. Then
\begin{equation}
  \sum_{j=1}^{g} \mathrm{MAC}(\psi, e_j) = \sum_{j=1}^{g} |\langle e_j, \psi\rangle|^2,
\end{equation}
which by Pythagoras is the squared length of the part of $\psi$ that lies inside the
group. It equals 1 exactly if $\psi$ lies entirely in the group, and it does not depend
on which basis the solver chose. The neighbouring mode sits at a slightly different
$\mathbf{k}$, so the sum is only approximately 1. Degenerate groups therefore have to
be compared as a whole, and all members of the group must have been computed.

\paragraph{3. A true crossing.} Two sorted bands swap identities: the sorted MAC is
about 0 and the MAC with the other band is about 1. In the sweep, bands 5 and 6 near
$k \approx 0.33$ on $M$--$\Gamma$ (about $6.5\,\mathrm{kHz}$, above the $5\,\mathrm{kHz}$
range) show a MAC of $0.0000$. That is consistent with two modes of opposite symmetry
about the diagonal, which have $c = 0$; the symmetry itself was not checked.

% ---------------------------------------------------------------------
\section{Tracking and gap finding are different tasks}

\emph{Tracking} follows a mode's identity from one $\mathbf{k}$ to the next, using the
MAC between neighbours. It answers ``which curve is this dot on?''

\emph{Gap finding} asks whether any dot lies in a frequency interval. It needs only the
frequencies sorted at each $\mathbf{k}$: the gap between sorted bands $n$ and $n+1$
runs from the maximum of band $n$ to the minimum of band $n+1$.

A tracked band may cross another one, so the $n$-th tracked band is not the $n$-th
sorted band. Suppose tracked curves $A$ and $B$ cross and a third band $C$ lies above
both. The sorted band below $C$ is the upper envelope of $A$ and $B$ at each
$\mathbf{k}$, so the correct lower gap edge is the maximum of that envelope. Taking the
curve $B$ alone when $A$ rises higher later on the path gives a value that is too low:
the reported gap starts too low and contains dots of $A$. Stage 5 therefore keeps plain
sorting available.

% ---------------------------------------------------------------------
\section{Check questions with answers}

\paragraph{Question 1.} \emph{Band 3 near $\Gamma$ has a neighbouring MAC of about 0.2,
and it stays there when the number of $\mathbf{k}$-points is doubled. What do you
conclude, and what would you do differently for that step?}

\emph{Answer.} Band 3 meets two other bands at $\Gamma$, and the solver returns an
arbitrary basis of that group. The MAC has to be summed over the whole group, which
then gives approximately 1. A degeneracy has to be told apart from a true crossing
between two sampled points: there the continuation is well defined, the MAC with the
partner band is about 1 and the sorted labels simply swap. Here the MAC is about 0.2,
neither 0 nor 1, and refining the step does not change it, which points to a degenerate
point where individual vectors mean nothing. The sum is only approximately 1 because the
neighbouring vector sits at a slightly different $\mathbf{k}$ (measured: $0.97$ to
$1.00$), and it needs every member of the group, which is why the solver uses guard
eigenpairs. For that step: match the group as a whole instead of one vector to one
vector. At the degenerate dot all members have the same frequency, so the plotted
curve does not depend on which label goes where.

\paragraph{Question 2.} \emph{Band 1 changes from $80\,\%$ plate to $5\,\%$ plate along
$M$--$\Gamma$, yet the neighbouring MAC stays close to 1 there. Why is this not a
contradiction, and which change in the unit cell would make this crossing harder to
track?}

\emph{Answer.} The change is gradual, spread over essentially all 57 steps of the
segment ($69\,\%$ to $24\,\%$ alone takes 12 steps), and each step only compares two
nearly identical shapes. A lighter resonator makes the crossing narrower, because the
coupling is $c = -\sqrt{\mu}\,\lambda_r$; the turn then happens within fewer steps, and
at the same step size the MAC dips further from 1.

\paragraph{Question 3.} \emph{Two tracked bands overlap in frequency, with a true crossing
between them. Could there still be a complete gap between sorted bands $n$ and $n+1$ at
that crossing frequency, and which data decides it?}

\emph{Answer.} No. At a true crossing the sorted bands $n$ and $n+1$ share a dot at the
crossing frequency $f_c$, so $\max(\text{band } n) \ge f_c \ge \min(\text{band } n+1)$
and there is no gap between that pair. The data to decide this is the sorted frequencies
at every $\mathbf{k}$; no eigenvectors are needed. The point of the question is the
reverse danger: using tracked bands for gap edges can report a gap that contains dots of
another band, as shown in the previous section.

\paragraph{Follow-up: how does the matrix $H$ of the two-level model translate into the
notation with mean, half difference and mixing angle?}
\emph{Answer.} This is worked out in Section~\ref{sec:mixing_angle}: $H_{11}$ and
$H_{22}$ are replaced by $\bar{H} \pm \Delta$, the off-diagonal $c$ stays the same, and
the eigenvectors depend only on $c/\Delta$.

% =====================================================================

\chapter{Eigenvectors, resonator participation and the MAC}
\label{ch:eigenvectors}

So far only the eigenvalues of the Bloch problem were used. This chapter looks at
the eigenvectors: what they contain, how much of a mode sits in the resonator, how
two modes at neighbouring wavevectors are compared with the Modal Assurance
Criterion (MAC), and which situations a mode tracker has to handle. All FE numbers
are measured on the production cell: $a = 0.1\,\mathrm{m}$, $h = 5\,\mathrm{mm}$,
steel, $10\times 10$ Lagrange degree 2 elements, one resonator at the cell centre
with $f_r = 1200\,\mathrm{Hz}$ and $\mu = 0.25$.

\begin{center}
\begin{tabular}{|l|p{0.68\textwidth}|}
\hline
Symbol & Meaning \\
\hline
$K_r$, $M_r$ & reduced stiffness and mass matrices at one $\mathbf{k}$, $K_r = T^H K T$, $M_r = T^H M T$ \\
$T(\mathbf{k})$ & maps independent DOFs to all DOFs, with phase factors \\
$\lambda = \omega^2$ & eigenvalue of one dot in the diagram \\
$\psi$ & reduced eigenvector of one dot, the column returned by the solver \\
$u$ & last entry of $\psi$: the resonator displacement \\
$m$, $k_s$ & resonator mass and spring stiffness \\
$\lambda_r = k_s/m$ & the resonator's own eigenvalue \\
$\lambda_p$ & bare plate eigenvalue of one plate wave \\
$\mu = m/(\rho h A)$ & resonator mass over plate mass per cell \\
$\langle \psi_1, \psi_2\rangle$ & mass-weighted inner product $\psi_1^H M_r \psi_2$ \\
$H$ & $2\times 2$ standard-form matrix of the two-level model \\
$q$, $v$ & entries of the eigenvector of $H$: plate wave amplitude and $v = \sqrt{m}\,u$ \\
$\bar{H}$, $\Delta$ & mean and half difference of the diagonal entries of $H$ \\
$c$ & off-diagonal entry of $H$ \\
$R$ & $\sqrt{\Delta^2 + c^2}$ \\
$\chi$ & mixing angle of the $2\times 2$ eigenvector \\
$k$ (in tables) & position on the diagonal, $k_x = k_y = k\,\pi/a$ \\
\hline
\end{tabular}
\end{center}

% ---------------------------------------------------------------------
\section{What an eigenvector is at one wavevector}

At each $\mathbf{k}$ the solver returns pairs $(\lambda, \psi)$ with
\begin{equation}
  K_r \psi = \lambda M_r \psi .
\end{equation}
Each dot in the diagram owns one $\psi$. Its entries are the plate unknowns
(deflection and two rotations) at the independent nodes, and its last entry is the
resonator displacement $u$. The resonator has no periodic partner, so $T$ acts on it
as the identity, and the last row of $M_r$ contains only $m$ on the diagonal.

$\psi$ is defined only up to a factor: if $\psi$ solves the equation, so does
$\gamma\psi$ for any complex number $\gamma$. Its scale and its overall complex
phase therefore carry no physical meaning, and everything computed from $\psi$ has
to be independent of both.

% ---------------------------------------------------------------------
\section{Orthogonality and what a degeneracy does}
\label{sec:orthogonality}

Take two solutions at the same $\mathbf{k}$:
\begin{equation}
  K_r \psi_1 = \lambda_1 M_r \psi_1, \qquad K_r \psi_2 = \lambda_2 M_r \psi_2 .
\end{equation}
Multiply the first from the left by $\psi_2^H$:
\begin{equation}
  \psi_2^H K_r \psi_1 = \lambda_1\, \psi_2^H M_r \psi_1 .
\end{equation}
Take the conjugate transpose of the second. $K_r$ and $M_r$ are Hermitian and
$\lambda_2$ is real, so $\psi_2^H K_r = \lambda_2\, \psi_2^H M_r$. Multiply from the
right by $\psi_1$:
\begin{equation}
  \psi_2^H K_r \psi_1 = \lambda_2\, \psi_2^H M_r \psi_1 .
\end{equation}
The left sides are equal. Subtracting gives
\begin{equation}
  (\lambda_1 - \lambda_2)\, \psi_2^H M_r \psi_1 = 0 .
\end{equation}
Different eigenvalues therefore mean $\psi_2^H M_r \psi_1 = 0$: the vectors are
orthogonal when weighted with $M_r$. That is why the mass matrix appears in
everything below.

If $\lambda_1 = \lambda_2 = \lambda$, the equation says nothing. Then every
combination is also a solution,
\begin{equation}
  K_r(\alpha_1\psi_1 + \alpha_2\psi_2) = \lambda M_r(\alpha_1\psi_1 + \alpha_2\psi_2),
\end{equation}
so at a degenerate point the solver returns \emph{some arbitrary basis} of the
group. The individual vectors of the group have no meaning of their own; only the
space they span does.

In this cell, degenerate groups appear where several bare plate waves share one
frequency. One combination couples to the resonator and leaves the group; the rest
do not see it. For example, $M$ has four plate waves at $2378\,\mathrm{Hz}$: one
coupled combination moves away and three modes remain at $2379\,\mathrm{Hz}$.

% ---------------------------------------------------------------------
\section{Plate share in the FE model}

The kinetic energy of a mode is $\tfrac12\omega^2\,\psi^H M_r\psi$, and the part in
the resonator is $\tfrac12\omega^2\, m|u|^2$. Hence
\begin{equation}
  \text{plate share} = 1 - \frac{m|u|^2}{\psi^H M_r \psi} .
\end{equation}
Scale and phase of $\psi$ cancel, as required.

\begin{center}
\begin{tabular}{|l|p{0.72\textwidth}|}
\hline
Point & Frequency in Hz: plate share (FE) \\
\hline
$\Gamma$ & $0$: $80\,\%$;\quad $1274$: $20\,\%$;\quad $4730$: $100\,\%$ (three modes);\quad $4883$: $100\,\%$ \\
$X$ & $836$: $31\,\%$;\quad $1195$: $100\,\%$;\quad $1653$: $69\,\%$;\quad $5881$: $100\,\%$ (three modes) \\
$M$ & $1031$: $5\,\%$;\quad $2379$: $100\,\%$ (three modes);\quad $2712$: $95\,\%$;\quad $11567$: $100\,\%$ \\
\hline
\end{tabular}
\end{center}

The simple model of Chapter~\ref{ch:local_resonance} gave $34\,\%$ and $66\,\%$ at
$X$; the shares at $\Gamma$ and $M$ agree with it.

\paragraph{The sum rule for the resonator share.}
Collect the eigenvectors of all bands at one $\mathbf{k}$ as columns of a matrix
$\Psi$, normalized so that $\Psi^H M_r \Psi = I$. Then $\Psi^{-1} = \Psi^H M_r$, so
\begin{equation}
  \Psi\,\Psi^H M_r = I \quad\Rightarrow\quad \Psi\,\Psi^H = M_r^{-1} .
\end{equation}
The resonator has no mass coupling to the plate, so the last diagonal entry of
$M_r^{-1}$ is $1/m$. The last diagonal entry of $\Psi\Psi^H$ is $\sum_n |u_n|^2$,
with $n$ running over all bands. Therefore
\begin{equation}
  \sum_n m|u_n|^2 = 1 .
\end{equation}
With this normalization $m|u_n|^2$ is the resonator share of band $n$, so
\emph{the resonator shares of all bands at one $\mathbf{k}$ add up to exactly
$100\,\%$}. The FE values at the three symmetry points show it: $20 + 80$ at
$\Gamma$, $69 + 31$ at $X$, $95 + 5$ at $M$. The plate shares of bands 1 and 2 alone
need not add up to $100\,\%$: along $M$--$\Gamma$ their sum drifts from $100\,\%$ to
about $104\,\%$ close to $M$, because other bands take part there.

% ---------------------------------------------------------------------
\section{The Modal Assurance Criterion}

\subsection{Definition and properties}
We need a number that says how alike two eigenvectors are and that ignores scale
and phase. With $\langle \psi_1, \psi_2\rangle = \psi_1^H M_r \psi_2$,
\begin{equation}
  \mathrm{MAC}(\psi_1, \psi_2) = \frac{|\langle\psi_1, \psi_2\rangle|^2}{\langle\psi_1, \psi_1\rangle\,\langle\psi_2, \psi_2\rangle} .
\end{equation}
\begin{itemize}
  \item Scale and phase drop out: replacing $\psi_1$ by $\gamma\psi_1$ multiplies
        numerator and denominator both by $|\gamma|^2$.
  \item By the Cauchy--Schwarz inequality it lies between 0 and 1. It is the squared
        cosine of the angle between the two vectors.
  \item $\mathrm{MAC} = 1$ for the same shape.
  \item $\mathrm{MAC} = 0$ for $M_r$-orthogonal vectors, which by
        Section~\ref{sec:orthogonality} holds for two different bands at the same
        $\mathbf{k}$.
\end{itemize}
Two vectors at neighbouring $\mathbf{k}$ come with two slightly different $M_r$. The
measurements below use $M_r$ of the first of the two points; the difference
vanishes as the step shrinks.

\subsection{How the step size enters}
Normalize $\langle\psi,\psi\rangle = 1$ and write the vector at the next
wavevector as
\begin{equation}
  \psi(\mathbf{k} + \delta\mathbf{k}) \approx \psi + \delta\,(\alpha\,\psi + p),
  \qquad \langle\psi, p\rangle = 0,
\end{equation}
with $\delta$ the step length, $\alpha$ a number and $p$ the part of the change
perpendicular to $\psi$. Then
\begin{align}
  \langle\psi, \psi(\mathbf{k} + \delta\mathbf{k})\rangle &= 1 + \delta\alpha, \\
  \langle\psi(\mathbf{k} + \delta\mathbf{k}), \psi(\mathbf{k} + \delta\mathbf{k})\rangle &= |1 + \delta\alpha|^2 + \delta^2\langle p, p\rangle ,
\end{align}
and so
\begin{equation}
  \mathrm{MAC} = \frac{|1 + \delta\alpha|^2}{|1 + \delta\alpha|^2 + \delta^2\langle p,p\rangle}
  \approx 1 - \delta^2\langle p, p\rangle .
\end{equation}
The part along $\psi$ only changes scale and phase and drops out. So
$1 - \mathrm{MAC} \propto \delta^2$, and halving the step should reduce it by a
factor of 4. Measured for band 1 at its worst step over the whole path (the point
number refers to the segments $\Gamma$--$X$ and $X$--$M$; $M$--$\Gamma$ gets
proportionally more):

\begin{center}
\begin{tabular}{|l|l|l|l|}
\hline
Points/segment & Wavevectors & Lowest MAC, band 1 & $1 - \mathrm{MAC}$ \\
\hline
20 & 69 & 0.9622 & 0.0378 \\
40 & 138 & 0.9898 & 0.0102 \\
80 & 274 & 0.9974 & 0.0026 \\
\hline
\end{tabular}
\end{center}

The ratios are $3.7$ and $3.9$, as expected. At 20 points per segment this worst
step lies just before $X$ on $\Gamma$--$X$, not at the resonator crossing on
$M$--$\Gamma$.

\subsection{Why only neighbours are compared}
Band 1 at $\Gamma$ against band 1 at $M$ has $\mathrm{MAC} = 0.35$, because the mode
legitimately changes from $80\,\%$ plate to $5\,\%$ plate. A low value over a long
distance means nothing. Between neighbouring points on $M$--$\Gamma$ the same band
never drops below $0.9974$ at 40 points per segment, even though it changes its
character completely along that path.

\subsection{Which vector and which weighting}
Lowest neighbouring MAC of band 1 at 40 points per segment:

\begin{center}
\begin{tabular}{|l|l|}
\hline
Variant & Lowest value \\
\hline
Full vector, $M_r$-weighted & 0.990 \\
Plate entries only, $M_r$-weighted & 0.969 \\
Full vector, unweighted & 0.619 \\
\hline
\end{tabular}
\end{center}

The MAC normalizes, so what is left of the plate motion in a resonator-dominated mode
still gives a usable comparison, only slightly worse. The clearly bad variant is
leaving out $M_r$: it counts every entry equally, whether deflection, rotation or
resonator, and falls to $0.62$ on a band that is perfectly smooth. Stage 5 therefore
uses the full vector including $u$, weighted with $M_r$.

% ---------------------------------------------------------------------
\section{How fast a vector turns at an avoided crossing}
\label{sec:mixing_angle}

This section takes the $2\times 2$ matrix of the two-level model apart, to see how
the eigenvector turns through an avoided crossing.

\subsection{The matrix we start from}
For one plate wave and one resonator, in mass-normalized coordinates, the eigenvector
is $(q, v)$ with $q$ the plate wave amplitude and $v = \sqrt{m}\,u$, and
Chapter~\ref{ch:local_resonance} gave
\begin{equation}
  H = \begin{pmatrix} H_{11} & c \\ c & H_{22} \end{pmatrix}
    = \begin{pmatrix} \lambda_p + \mu\lambda_r & -\sqrt{\mu}\,\lambda_r \\ -\sqrt{\mu}\,\lambda_r & \lambda_r \end{pmatrix} .
\end{equation}

\subsection{Mean and half difference}
Define the mean and the half difference of the diagonal,
\begin{equation}
  \bar{H} = \frac{H_{11} + H_{22}}{2}, \qquad \Delta = \frac{H_{11} - H_{22}}{2} .
\end{equation}
Adding and subtracting these gives $H_{11} = \bar{H} + \Delta$ and
$H_{22} = \bar{H} - \Delta$, so
\begin{equation}
  H = \bar{H}\, I + B, \qquad B = \begin{pmatrix} \Delta & c \\ c & -\Delta \end{pmatrix} .
\end{equation}
Inserting the entries:

\begin{center}
\begin{tabular}{|l|l|l|}
\hline
Symbol & In terms of $H$ & For plate wave and resonator \\
\hline
$\bar{H}$ & $(H_{11} + H_{22})/2$ & $\tfrac12\left[\lambda_p + \lambda_r(1+\mu)\right]$ \\
$\Delta$ & $(H_{11} - H_{22})/2$ & $\tfrac12\left[\lambda_p - \lambda_r(1-\mu)\right]$ \\
$c$ & off-diagonal entry & $-\sqrt{\mu}\,\lambda_r$ \\
\hline
\end{tabular}
\end{center}

\paragraph{Why split it.} If $Bz = rz$ for a vector $z$, then
$Hz = \bar{H}z + Bz = (\bar{H} + r)z$. The part $\bar{H}I$ only shifts the
eigenvalues. \emph{The eigenvectors depend on $B$ alone}, so only $\Delta$ and $c$
matter. $\Delta = 0$ at $\lambda_p = \lambda_r(1-\mu)$, where the diagonal entries
are equal: the centre of the avoided crossing.

\subsection{Eigenvalues}
\begin{align}
  \det(B - rI) &= (\Delta - r)(-\Delta - r) - c^2 = r^2 - \Delta^2 - c^2 = 0, \\
  \Rightarrow\quad r &= \pm R, \qquad R = \sqrt{\Delta^2 + c^2} .
\end{align}
The eigenvalues of $H$ are $\bar{H} \pm R$. This is the familiar two-level formula
\begin{equation}
  \lambda = \frac{H_{11} + H_{22}}{2} \pm \sqrt{\left(\frac{H_{11} - H_{22}}{2}\right)^2 + |c|^2} .
\end{equation}
Their sum is $2\bar{H} = \lambda_p + \lambda_r(1+\mu)$ and their product is
\begin{equation}
  \bar{H}^2 - R^2 = \bar{H}^2 - \Delta^2 - c^2 = H_{11}H_{22} - c^2
  = (\lambda_p + \mu\lambda_r)\lambda_r - \mu\lambda_r^2 = \lambda_p\lambda_r .
\end{equation}

\subsection{Eigenvectors and the mixing angle}
$(\Delta, c)$ is a point in a plane at distance $R$ from the origin. Call its polar
angle $2\chi$:
\begin{equation}
  \Delta = R\cos 2\chi, \qquad c = R\sin 2\chi, \qquad \tan 2\chi = \frac{c}{\Delta} .
\end{equation}
This is always possible because $\Delta^2 + c^2 = R^2$. Then
\begin{equation}
  B = R \begin{pmatrix} \cos 2\chi & \sin 2\chi \\ \sin 2\chi & -\cos 2\chi \end{pmatrix} .
\end{equation}
Apply $B$ to $(\cos\chi, \sin\chi)$ and use the addition theorems for
$\cos(2\chi - \chi)$ and $\sin(2\chi - \chi)$:
\begin{align}
  \text{row 1:} \quad & R(\cos 2\chi\cos\chi + \sin 2\chi\sin\chi) = R\cos(2\chi - \chi) = R\cos\chi, \\
  \text{row 2:} \quad & R(\sin 2\chi\cos\chi - \cos 2\chi\sin\chi) = R\sin(2\chi - \chi) = R\sin\chi .
\end{align}
So $(\cos\chi, \sin\chi)$ belongs to $+R$. Applying $B$ to $(-\sin\chi, \cos\chi)$
in the same way:
\begin{align}
  \text{row 1:} \quad & R(-\cos 2\chi\sin\chi + \sin 2\chi\cos\chi) = R\sin\chi = -R\,(-\sin\chi), \\
  \text{row 2:} \quad & R(-\sin 2\chi\sin\chi - \cos 2\chi\cos\chi) = -R\cos\chi ,
\end{align}
so it belongs to $-R$. In summary,
\begin{equation}
  \text{upper band: } (q, v) = (\cos\chi, \sin\chi), \qquad
  \text{lower band: } (q, v) = (-\sin\chi, \cos\chi) .
\end{equation}
The plate share is $|q|^2$. With $\cos^2\chi = (1 + \cos 2\chi)/2$ and
$\cos 2\chi = \Delta/R$,
\begin{equation}
  \boxed{\;\begin{aligned}
  \text{upper band: plate share} &= \tfrac12\left(1 + \frac{\Delta}{R}\right), \\
  \text{lower band: plate share} &= \tfrac12\left(1 - \frac{\Delta}{R}\right)
  \end{aligned}\;}
\end{equation}
Because $c < 0$ here, $\chi$ is negative. That only sets the relative sign of $q$ and
$v$; the shares do not depend on it.

\subsection{One worked point: $\Gamma$}
At $\Gamma$, $\lambda_p = 0$, and with $\mu = 0.25$:
\begin{align}
  H_{11} &= 0.25\,\lambda_r, \quad H_{22} = \lambda_r, \quad c = -0.5\,\lambda_r, \\
  \bar{H} &= 0.625\,\lambda_r, \quad \Delta = -0.375\,\lambda_r, \\
  R &= \sqrt{0.1406 + 0.25}\,\lambda_r = 0.625\,\lambda_r .
\end{align}
The eigenvalues are $\bar{H} \pm R = 0$ and $1.25\,\lambda_r$, the known $0$ and
$\lambda_r(1+\mu)$. With $\Delta/R = -0.6$, the lower band has plate share
$\tfrac12(1 + 0.6) = 80\,\%$ and the upper band $\tfrac12(1 - 0.6) = 20\,\%$.
Further, $2\chi = -126.9^\circ$, so $\chi = -63.4^\circ$.

This agrees with the ratio formula of Chapter~\ref{ch:local_resonance},
$q/v = (\lambda_r - \lambda)/(\sqrt{\mu}\,\lambda_r)$: for the upper band it gives
$-0.25/0.5 = -0.5$, and $\cot(-63.4^\circ) = -0.5$.

It also explains why band 1 starts at $80\,\%$ and not at $100\,\%$. Since $\lambda_p$
cannot go below 0, $\Delta$ cannot go below $-0.375\,\lambda_r$, and the far end of
the turn is never reached.

\subsection{Along $M$--$\Gamma$: model against FE}
Plate shares in percent, for bands 1 and 2 (B1, B2), one-wave model (mod.)\ against
the FE cell:

\begin{center}
\begin{tabular}{|l|l|l|l|l|l|l|l|}
\hline
$k$ & $\lambda_p/\lambda_r$ & $\Delta/\lambda_r$ & $\chi$ & B1 mod. & B1 FE & B2 mod. & B2 FE \\
\hline
0.018 & 0.000 & $-0.375$ & $-63.4^\circ$ & 80.0 & 80.0 & 20.0 & 20.4 \\
0.509 & 0.267 & $-0.242$ & $-57.9^\circ$ & 71.8 & 70.3 & 28.2 & 31.2 \\
0.649 & 0.705 & $-0.023$ & $-46.3^\circ$ & 52.3 & 46.3 & 47.7 & 56.1 \\
0.719 & 1.060 & $0.155$ & $-36.4^\circ$ & 35.2 & 27.7 & 64.8 & 75.2 \\
0.789 & 1.536 & $0.393$ & $-25.9^\circ$ & 19.1 & 14.2 & 80.9 & 89.2 \\
0.930 & 2.943 & $1.096$ & $-12.3^\circ$ & 4.5 & 5.6 & 95.5 & 98.9 \\
1.000 & 3.928 & $1.589$ & $-8.7^\circ$ & 2.3 & 5.1 & 97.7 & 100.0 \\
\hline
\end{tabular}
\end{center}

The model has $\Delta = 0$ at $k \approx 0.667$, and the FE band 1 passes $50\,\%$
shortly before $k = 0.649$. The agreement fades towards $M$, as it should: at $M$
four plate waves take part, and this matrix holds only one. The band frequencies show
the same (band 1 at $M$: $1153\,\mathrm{Hz}$ in the model, $1031\,\mathrm{Hz}$ in
FE). With these values $\chi$ runs from $-63.4^\circ$ at $\Gamma$ to $-8.7^\circ$ at
$M$, about $55^\circ$ in total, not the idealized full turn of $90^\circ$.

\subsection{From the angle to the MAC}
Two normalized vectors at angles $\chi_1$ and $\chi_2$ have the inner product
$\cos\chi_1\cos\chi_2 + \sin\chi_1\sin\chi_2 = \cos(\chi_1 - \chi_2)$, so
\begin{equation}
  \mathrm{MAC} = \cos^2(\chi_1 - \chi_2) .
\end{equation}
To see how fast $\chi$ turns, differentiate $2\chi = \arctan(c/\Delta)$ with respect
to $\Delta$:
\begin{equation}
  2\,\frac{\mathrm{d}\chi}{\mathrm{d}\Delta}
  = \frac{1}{1 + c^2/\Delta^2}\cdot\left(-\frac{c}{\Delta^2}\right)
  = -\frac{c}{\Delta^2 + c^2}
  = -\frac{c}{R^2} .
\end{equation}
The turn is fastest at $\Delta = 0$, where $R = |c|$ and
$|\mathrm{d}\chi/\mathrm{d}\Delta| = 1/(2|c|)$. Since $\Delta$ contains $\lambda_p/2$,
\begin{equation}
  \left|\frac{\mathrm{d}\chi}{\mathrm{d}\lambda_p}\right|_{\Delta = 0}
  = \frac{1}{4|c|} = \frac{1}{4\sqrt{\mu}\,\lambda_r} .
\end{equation}
Half of the full $90^\circ$ turn happens while $|\Delta| < |c|$: at $\Delta = |c|$ one
has $2\chi = -45^\circ$, at $\Delta = -|c|$ one has $2\chi = -135^\circ$, so $\chi$
changes by $45^\circ$. In terms of $\lambda_p$ this window has the width
$4\sqrt{\mu}\,\lambda_r$: $2.0\,\lambda_r$ for $\mu = 0.25$, $0.4\,\lambda_r$ for
$\mu = 0.01$. A lighter resonator has a narrower crossing and needs smaller steps.

Measured for band 1 on $M$--$\Gamma$, with the spacing $0.0175\,\pi/a$ of the 40-point
path (57 steps on this segment):

\begin{center}
\begin{tabular}{|l|l|l|}
\hline
 & Lowest MAC on $M$--$\Gamma$ & at $k$ \\
\hline
Model, $\cos^2(\chi_1 - \chi_2)$ & 0.99778 & 0.737 \\
FE & 0.99736 & 0.702 \\
\hline
\end{tabular}
\end{center}

The minimum lies a little past $\Delta = 0$ because $\lambda_p$ also changes faster
with $k$ as $k$ grows. The change from $69\,\%$ to $24\,\%$ plate takes 12 steps, and
the change from $80\,\%$ to $5\,\%$ takes practically all 57.

\paragraph{The baseline of the FE value.}
Close to $\Gamma$ the FE MAC stays at $0.99961$ while the model gives $1.00000$. The
model uses ``the plate wave at this $\mathbf{k}$'' as a basis vector, so it cannot see
that the pattern $e^{i\mathbf{k}\cdot\mathbf{x}}$ itself changes between neighbouring
points. One step changes the phase at $\mathbf{x}$ by $\theta(\mathbf{x}) =
\delta\mathbf{k}\cdot\mathbf{x}$, and $1 - \mathrm{MAC}$ is approximately the variance
of $\theta$ over the mass of the mode. With $\delta k_x = \delta k_y = 0.0175\,\pi/a$,
a variance of $a^2/12$ of each coordinate over the cell, and $80\,\%$ of the mass
motion in the plate (the resonator sits at the mean position):
\begin{equation}
  1 - \mathrm{MAC} \approx 0.8\cdot (0.0175\,\pi)^2\cdot\frac{2}{12} \approx 0.0004 ,
\end{equation}
which matches the measured $1 - 0.99961 = 0.00039$.

\subsection{The other two-level matrix: two plate waves and a rigid mass}
Chapter~\ref{ch:two_level} used the same structure with
$H_{11} = \lambda_1^0 - \mu\lambda$, $H_{22} = \lambda_2^0 - \mu\lambda$ and
$c = -\mu\lambda$. Then
\begin{equation}
  \bar{H} = \frac{\lambda_1^0 + \lambda_2^0}{2} - \mu\lambda, \qquad
  \Delta = \frac{\lambda_1^0 - \lambda_2^0}{2},
\end{equation}
because the mass loading cancels in the difference. At $X$, $\Delta = 0$, so
$2\chi = -90^\circ$ and $\chi = -45^\circ$. That gives the eigenvectors
$(1, -1)/\sqrt{2}$ and $(1, 1)/\sqrt{2}$ of that chapter, with $a_H = \bar{H}$ and
$b_H = c$.

% ---------------------------------------------------------------------
\section{The three situations a tracker meets}

Lowest neighbouring MAC between \emph{sorted} bands along the whole path:

\begin{center}
\begin{tabular}{|l|l|l|l|p{0.32\textwidth}|}
\hline
Band & 20 points & 40 points & 80 points & Where \\
\hline
1 & 0.9622 & 0.9898 & 0.9974 & smooth band \\
2 & 0.2301 & 0.2325 & 0.2322 & next to $M$, three modes at $2379\,\mathrm{Hz}$ \\
3 & 0.1755 & 0.1954 & 0.2117 & next to $\Gamma$, three modes at $4730\,\mathrm{Hz}$ \\
4 & 0.0601 & 0.0610 & 0.0615 & next to $X$, three modes at $5881\,\mathrm{Hz}$ \\
5 & 0.0000 & 0.0000 & 0.0000 & near $k \approx 0.33$ on $M$--$\Gamma$, about $6.5\,\mathrm{kHz}$ \\
\hline
\end{tabular}
\end{center}

\paragraph{1. Smooth continuation.} Band 1: the MAC is close to 1 and improves with a
finer step.

\paragraph{2. A degenerate point.} Bands 2 to 4: the value barely changes with a finer
step. That is the signature of a degeneracy and not of a coarse step. At the
degenerate point the solver picks an arbitrary basis of the group
(Section~\ref{sec:orthogonality}), while just beside it the modes are distinct. Near
$M$, MAC of the three neighbouring modes against the three modes at $M$ (40 points per
segment):

\begin{center}
\begin{tabular}{|l|l|l|l|l|}
\hline
Before $M$ & $M$, mode 1 & $M$, mode 2 & $M$, mode 3 & Row sum \\
\hline
$2321\,\mathrm{Hz}$ & 0.497 & 0.045 & 0.457 & 1.000 \\
$2372\,\mathrm{Hz}$ & 0.017 & 0.882 & 0.072 & 0.971 \\
$2439\,\mathrm{Hz}$ & 0.485 & 0.047 & 0.468 & 1.000 \\
\hline
\end{tabular}
\end{center}

No single entry tells the continuation, but the row sums are close to 1. The reason:
let $e_1, \dots, e_g$ be an $M_r$-orthonormal basis of the group and $\psi$ a vector
with $\langle\psi,\psi\rangle = 1$. Then
\begin{equation}
  \sum_{j=1}^{g} \mathrm{MAC}(\psi, e_j) = \sum_{j=1}^{g} |\langle e_j, \psi\rangle|^2,
\end{equation}
which by Pythagoras is the squared length of the part of $\psi$ that lies inside the
group. It equals 1 exactly if $\psi$ lies entirely in the group, and it does not depend
on which basis the solver chose. The neighbouring mode sits at a slightly different
$\mathbf{k}$, so the sum is only approximately 1. Degenerate groups therefore have to
be compared as a whole, and all members of the group must have been computed.

\paragraph{3. A true crossing.} Two sorted bands swap identities: the sorted MAC is
about 0 and the MAC with the other band is about 1. In the sweep, bands 5 and 6 near
$k \approx 0.33$ on $M$--$\Gamma$ (about $6.5\,\mathrm{kHz}$, above the $5\,\mathrm{kHz}$
range) show a MAC of $0.0000$. That is consistent with two modes of opposite symmetry
about the diagonal, which have $c = 0$; the symmetry itself was not checked.

% ---------------------------------------------------------------------
\section{Tracking and gap finding are different tasks}

\emph{Tracking} follows a mode's identity from one $\mathbf{k}$ to the next, using the
MAC between neighbours. It answers ``which curve is this dot on?''

\emph{Gap finding} asks whether any dot lies in a frequency interval. It needs only the
frequencies sorted at each $\mathbf{k}$: the gap between sorted bands $n$ and $n+1$
runs from the maximum of band $n$ to the minimum of band $n+1$.

A tracked band may cross another one, so the $n$-th tracked band is not the $n$-th
sorted band. Suppose tracked curves $A$ and $B$ cross and a third band $C$ lies above
both. The sorted band below $C$ is the upper envelope of $A$ and $B$ at each
$\mathbf{k}$, so the correct lower gap edge is the maximum of that envelope. Taking the
curve $B$ alone when $A$ rises higher later on the path gives a value that is too low:
the reported gap starts too low and contains dots of $A$. Stage 5 therefore keeps plain
sorting available.

% ---------------------------------------------------------------------
\section{Check questions with answers}

\paragraph{Question 1.} \emph{Band 3 near $\Gamma$ has a neighbouring MAC of about 0.2,
and it stays there when the number of $\mathbf{k}$-points is doubled. What do you
conclude, and what would you do differently for that step?}

\emph{Answer.} Band 3 meets two other bands at $\Gamma$, and the solver returns an
arbitrary basis of that group. The MAC has to be summed over the whole group, which
then gives approximately 1. A degeneracy has to be told apart from a true crossing
between two sampled points: there the continuation is well defined, the MAC with the
partner band is about 1 and the sorted labels simply swap. Here the MAC is about 0.2,
neither 0 nor 1, and refining the step does not change it, which points to a degenerate
point where individual vectors mean nothing. The sum is only approximately 1 because the
neighbouring vector sits at a slightly different $\mathbf{k}$ (measured: $0.97$ to
$1.00$), and it needs every member of the group, which is why the solver uses guard
eigenpairs. For that step: match the group as a whole instead of one vector to one
vector. At the degenerate dot all members have the same frequency, so the plotted
curve does not depend on which label goes where.

\paragraph{Question 2.} \emph{Band 1 changes from $80\,\%$ plate to $5\,\%$ plate along
$M$--$\Gamma$, yet the neighbouring MAC stays close to 1 there. Why is this not a
contradiction, and which change in the unit cell would make this crossing harder to
track?}

\emph{Answer.} The change is gradual, spread over essentially all 57 steps of the
segment ($69\,\%$ to $24\,\%$ alone takes 12 steps), and each step only compares two
nearly identical shapes. A lighter resonator makes the crossing narrower, because the
coupling is $c = -\sqrt{\mu}\,\lambda_r$; the turn then happens within fewer steps, and
at the same step size the MAC dips further from 1.

\paragraph{Question 3.} \emph{Two tracked bands overlap in frequency, with a true crossing
between them. Could there still be a complete gap between sorted bands $n$ and $n+1$ at
that crossing frequency, and which data decides it?}

\emph{Answer.} No. At a true crossing the sorted bands $n$ and $n+1$ share a dot at the
crossing frequency $f_c$, so $\max(\text{band } n) \ge f_c \ge \min(\text{band } n+1)$
and there is no gap between that pair. The data to decide this is the sorted frequencies
at every $\mathbf{k}$; no eigenvectors are needed. The point of the question is the
reverse danger: using tracked bands for gap edges can report a gap that contains dots of
another band, as shown in the previous section.

\paragraph{Follow-up: how does the matrix $H$ of the two-level model translate into the
notation with mean, half difference and mixing angle?}
\emph{Answer.} This is worked out in Section~\ref{sec:mixing_angle}: $H_{11}$ and
$H_{22}$ are replaced by $\bar{H} \pm \Delta$, the off-diagonal $c$ stays the same, and
the eigenvectors depend only on $c/\Delta$.
